\documentclass[12pt]{article}
\input{preamble}

\begin{document}
\pagenumbering{Roman}

\begin{titlepage}
  \centering
  \vspace*{3cm}
  {\Huge\bfseries Agentic Deep Research System\par}
  \vspace{0.8cm}
  {\LARGE 面试技术拷打笔记（TeX 正式版）\par}
  \vspace{1.2cm}
  {\large 面向：技术面试 / 架构深挖 / 治理闭环追问 / 高压交叉盘问\par}
  \vfill
  {\large 版本：v2.1\par}
  {\large 日期：2026-07-27\par}
\end{titlepage}

\clearpage
\setcounter{page}{2}
\pdfbookmark[1]{目录}{toc}
\tableofcontents
\clearpage
\pagenumbering{arabic}

\section{怎么使用这份笔记}

\begin{summaryBox}{这份文档不是项目报告}
这份文档的用途不是把项目按模块介绍一遍，而是训练一种更适合技术面试的表达方式：面试官会先确认你是不是自己做过，再追问你为什么这么选，接着看你能不能把边界、缺陷、治理和演进路线讲清楚。
\end{summaryBox}

\begin{qaBox}{建议的回答顺序}
回答时不要一上来就堆框架名。最稳的顺序是：先定义系统解决什么问题，再讲主链路怎么跑，再讲为什么这么选，再讲当前已经做到哪里，最后讲哪些地方还没做完、准备怎么补。这样既像工程师，也不像背稿。
\end{qaBox}

\begin{analysisBox}{使用方法}
这份笔记更适合按“结论、深入分析、边界”三层来准备。先用一两句话给出立场，再补业务抽象、替代方案、代价和故障模式，最后明确哪些能力已经落地、哪些仍有边界。面试官真正想听的通常不是定义本身，而是你是否能在压力下把判断讲完整。
\end{analysisBox}

\begin{riskBox}{这份笔记的使用原则}
第一，不要把设计稿说成已上线功能。第二，不要把历史 bug 说成当前仍存在的问题。第三，不要把治理愿景说成生产级闭环。第四，遇到面试官追问时，优先讲事实、边界和取舍，不要用空话兜底。
\end{riskBox}

\section{第一轮拷打：先把项目说清楚}

\begin{summaryBox}{这一轮面试官在试什么}
这一轮不是考你会不会背名词，而是在判断你能不能用两到三分钟把系统定义准确。面试官真正想听的是：这到底是不是一个严肃的研究型系统，还是只是把搜索、RAG 和报告拼到一起的 Demo。
\end{summaryBox}

\begin{qaBox}{这一轮最常见的追问}
\begin{itemize}[leftmargin=2em]
  \item 这个项目一句话到底是什么，不要给我讲一堆模块名。
  \item 它和普通 RAG、普通搜索问答、普通工作流引擎到底差在哪。
  \item 你这个系统的输入、处理中间态、输出结果分别是什么。
  \item 如果我只给你两分钟，你最希望我记住这个项目的哪三个关键词。
\end{itemize}
\end{qaBox}

\subsection{先用一句话定义这个项目}
\begin{qaBox}{一句话总述}
这是一个以 LangGraph 为业务编排内核、以 Search / Browser / RAG 为取证能力、以 Analyst / Reflection / Report 为分析闭环、并逐步向治理型 Agent 演进的研究系统。
\end{qaBox}

\subsection{面试开场应该怎么讲}
\begin{summaryBox}{推荐开场}
如果面试官先问“这个系统是干什么的”，最稳的讲法不是直接说“这是一个 AI Agent”，而是先把它定义成一个 \term{evidence-driven deep research workflow}。它的核心目标不是闲聊，不是通用办公自动化，而是围绕“搜索、取证、对比、分析、引用、报告”构造一条可重复执行的研究链路。
\end{summaryBox}

我通常会用下面这段话开场：

\begin{qaBox}{标准开场答案}
这个系统的目标是把用户问题转换成一张可执行研究图。Planner 先把问题拆成 DAG，Search / Browser / RAG 节点分别去互联网、网页和内部知识库取证，Analyst 负责综合证据形成结构化分析，Reflection 负责做事实、一致性和引用覆盖率校验，最后 Report 输出带引用的研究报告。系统的设计重点不是单次回答得多花哨，而是证据链、失败语义、预算控制和后续治理扩展能力。
\end{qaBox}

\subsection{系统的分层怎么看}
\begin{itemize}[leftmargin=2em]
  \item 编排层：LangGraph StateGraph，负责节点、条件边、循环和并行。
  \item 能力层：Planner / Search / Browser / RAG / Analyst / Reflection / Report。
  \item 状态层：ResearchState、DAG、checkpoint、trace、tool audit、budget state。
  \item 数据层：PostgreSQL + pgvector，存文档、session、citation、tool audit、budget state。
  \item 交互层：FastAPI + SSE，负责创建任务、查询状态、流式推送进度。
  \item 治理层：guardrail、runtime status、budget breaker、approval source of truth、Harness supervisor、MCP proxy。
\end{itemize}

\subsection{为什么这个系统不是简单的 RAG 问答}
\begin{qaBox}{核心差异}
RAG 问答通常是“检索一轮，再回答一轮”。这个系统做的是“先规划，再多源取证，再综合分析，再反思校验，再决定是否重规划”。它把单次回答升级成了一条带状态、带条件分支、带恢复与治理约束的研究流水线。
\end{qaBox}

\begin{analysisBox}{深入分析}
这道题不能只停在“步骤更多”。面试官真正想确认的是，你做的到底是一个回答器，还是一个研究执行系统。普通 RAG 的默认假设是：问题边界比较清楚，检索一次拿到上下文后，模型就可以直接组织答案。所以它的核心矛盾是“召回准不准、答案顺不顺”。而研究型 Agent 的假设完全不同，它默认问题可能被拆错、证据可能互相冲突、某个来源可能暂时不可用、预算可能不允许无限扩张、以及第一次答案可能需要被系统自己否掉重做。

因此两者的差异不只是模块数量，而是执行模型。简单 RAG 的主对象是“上下文窗口”，研究型系统的主对象是“任务状态机”。前者关注把哪些文档塞进 prompt，后者关注当前任务处于哪一阶段、还有哪些节点没跑、失败是可重试还是终止、是否允许继续联网、是否需要人工审批、是否还有预算继续扩张。只要把主对象换掉，系统设计就会从“检索增强回答”自然演进成“可治理的研究流程”。

这一点在面试里必须讲透，因为很多人会把 Search、Browser、RAG、Report 都堆在一起，自称做了 Agent。真正的分水岭不是工具数量，而是是否有显式状态、显式路由、显式失败语义和显式治理约束。你要把这层抽象抬出来，面试官才会认为你不是在做功能拼装。
\end{analysisBox}

\section{第二轮拷打：为什么这么选，而不是别的}

\begin{summaryBox}{这一轮面试官在试什么}
当你把系统讲清楚之后，面试官通常会立刻追问选型理由。这里最怕两种答法：一种是“社区都这么用”，另一种是“我觉得这个更先进”。正确答法必须同时包含备选方案、比较维度和最终取舍。
\end{summaryBox}

\begin{riskBox}{这一轮最容易翻车的地方}
不要把“能用”讲成“唯一正确”。也不要回避被否选方案，因为真正做过工程的人一定经历过权衡，而不是从一开始就知道唯一答案。
\end{riskBox}

\begin{qaBox}{这一轮最常见的追问}
\begin{itemize}[leftmargin=2em]
  \item 为什么是 LangGraph，不是 LangChain、AutoGen、CrewAI，或者你自己写调度器。
  \item 为什么浏览器自动化选 Playwright，不是 Puppeteer 或 Selenium。
  \item 为什么检索层放 pgvector，不直接用专门向量库。
  \item 为什么还要 Harness，说明 LangGraph 还不够，那你一开始为什么还选它。
\end{itemize}
\end{qaBox}

\subsection{Agent 编排框架：为什么选择 LangGraph}
\begin{qaBox}{结论}
研究任务本质上不是线性 \term{Chain}，而是带并行、条件分支和重规划循环的状态机，所以选择 LangGraph，而不是只用 LangChain 的串联式抽象。
\end{qaBox}

\begin{center}
\small
\begin{tabularx}{\textwidth}{>{\bfseries}lXXXX}
\toprule
维度 & LangGraph & LangChain & AutoGen & CrewAI \\
\midrule
核心抽象 & 状态机 / 图编排 & 链和工具集成层 & 对话式多 Agent & 角色协作脚本 \\
并行与 fan-out & 原生 \term{Send API} & 需手动拼接 & 可做但偏消息流 & 较弱 \\
条件边 & 强 & 中 & 中 & 弱 \\
循环与重规划 & 强 & 弱 & 中 & 弱 \\
检查点恢复 & 支持 & 有限 & 有限 & 弱 \\
适合研究 DAG & 很适合 & 一般 & 偏实验 & 偏演示 \\
\bottomrule
\end{tabularx}
\end{center}

\begin{summaryBox}{选择理由}
这里至少有两个可行方案。方案 A 是继续用 LangChain，把 Planner、Search、Report 串起来，再自己维护一套 if-else。优点是简单，缺点是并行、循环、恢复都会越来越别扭。方案 B 是直接把研究流程建模成一张图，让节点、边、状态、路由都显式化。最终选择 LangGraph，就是选择方案 B。
\end{summaryBox}

\begin{qaBox}{面试回答模板}
LangGraph 解决的是“复杂多步骤任务如何以显式状态机方式执行”。研究型任务天然需要 DAG、并行、条件边和循环重规划，这四点是 LangGraph 的强项。LangChain 很适合做工具和模型的胶水层，但不适合单独承担复杂业务工作流编排。
\end{qaBox}

\begin{analysisBox}{深入分析}
这一题真正的考点不是“你知不知道 LangGraph”，而是“你有没有把业务问题正确映射到执行模型上”。如果业务天然是线性的，比如固定三步：检索、总结、输出，那么 LangChain 的链式抽象完全够用，甚至更轻。但研究型任务不一样，它会同时出现三类复杂性。第一类是并行性，你希望 search、browser、rag 在某些阶段同时跑，而不是人为串起来浪费 I/O 时间。第二类是条件性，你要根据 RAG 命中情况决定是否继续联网，根据 Reflection 结果决定是否 replan，根据审批结果决定是否恢复。第三类是循环性，你不能假设第一次规划一定正确，所以必须给重规划留合法路径。

如果用普通 chain 硬做这些事，短期能跑，长期会出现两个工程问题。一个问题是状态散落。你会把当前步骤、已完成节点、预算、失败次数、审批状态藏在 if-else 和局部变量里，导致一旦出 bug 很难解释“系统为什么走到了这里”。另一个问题是恢复困难。因为链式抽象默认关注的是一步接一步的推理，而不是一个可中断、可恢复、可审计的长期工作流。

所以这里选择 LangGraph，本质上不是“选一个更潮的框架”，而是承认这个系统从第一天起就是状态机问题，而不是提示词编排问题。面试时如果你能把这层抽象对齐讲出来，可信度会比单纯报框架名高很多。
\end{analysisBox}

\subsection{LangGraph 和 DAG 的关系}
\begin{qaBox}{追问答案}
DAG 是业务层面的研究计划图，StateGraph 是执行层面的工作流图。Planner 生成的 DAG 解决“研究步骤怎么拆”，LangGraph StateGraph 解决“这些步骤怎样被调度、聚合、校验、循环和恢复”。DAG 是数据，StateGraph 是执行器。
\end{qaBox}

\begin{analysisBox}{深入分析}
这里最容易犯的错，是把 DAG 和 LangGraph 讲成同一个东西。实际上它们回答的是两个不同问题。DAG 是“这次研究任务需要做哪些事、谁依赖谁”，本质上是业务计划；LangGraph StateGraph 是“这些计划节点在运行时怎么被驱动、怎么聚合结果、怎么决定下一步”，本质上是执行引擎。一个是业务数据结构，一个是控制流框架。

为什么要分开？因为业务计划和执行控制的变化频率不同。Planner 今天可能多生成一种节点类型，或者把 Web 节点拆得更细，这属于 DAG 层的变化；但超时重试、审批暂停、反思重规划、预算熔断这类机制，不应该每次都由 Planner 直接负担，否则模型输出一旦不稳定，执行层就会跟着失控。把 DAG 当数据、把 StateGraph 当执行器，实际上是在保护控制面不被 LLM 计划输出直接污染。

还有一个关键收益是可解释性。如果 DAG 是数据，你就能清晰地把“计划是否合理”和“执行是否正确”拆开分析。前者看节点设计、依赖关系、fallback 是否健全；后者看路由、失败语义、恢复、预算和审计是否可靠。面试里把这两个层次拆开，会显得你对系统建模有清晰边界感。
\end{analysisBox}

\subsection{为什么还要 Harness，而不是只靠 LangGraph}
\begin{center}
\small
\begin{tabularx}{\textwidth}{>{\bfseries}lXXX}
\toprule
方案 & 做法 & 优点 & 缺点 \\
\midrule
方案 A & 生命周期、重试、审批、恢复都塞进 LangGraph & 单层编排、概念简单 & 图会越来越重，治理逻辑和业务逻辑耦合 \\
方案 B & LangGraph 管业务图，Harness 管 supervisor & 治理边界清晰，适合长任务 & 多一层状态同步与恢复控制 \\
\bottomrule
\end{tabularx}
\end{center}

\begin{qaBox}{最终选择}
当前项目的长期正确方向是方案 B：LangGraph 负责业务编排，Harness 负责任务级控制面。因为你真正要治理的是“失败后怎么恢复、预算超了怎么刹车、审批后怎么恢复、worker 怎么抢租约”，这些都不是业务 DAG 最擅长表达的东西。
\end{qaBox}

\begin{analysisBox}{深入分析}
这一题最容易答浅。很多人会说“LangGraph 不够，所以再加 Harness”，这句话不够严谨。更准确的说法是：LangGraph 擅长描述业务执行图，但不擅长承担完整的任务级控制面。业务执行图关心的是当前任务应该怎么做，supervisor 关心的是当前任务该不该继续做、由谁做、从哪里恢复做、预算是否允许做、审批是否允许做。两者不是强弱关系，而是职责分层关系。

如果把 supervisor 语义硬塞进图里，短期看起来只有一层系统，长期会有三个副作用。第一，图的节点语义会被污染。原本 analyst、reflection、report 这些节点是在表达业务，而不是在表达租约、暂停、恢复、超时、人工确认。第二，图会变得很难演进。因为每加一个治理策略，你都要再给图加分支、加状态、加条件，最后图就不再是业务图，而是一个混杂业务与控制的巨型控制器。第三，恢复语义会失真。任务级恢复往往需要依赖外部状态源，例如 harness 文件、worker 租约、checkpoint 版本，而这些都不是图内局部状态能优雅表达的。

所以外置 Harness 的关键价值在于分域定权：图内负责“怎么执行”，图外负责“是否继续执行、如何恢复执行”。这一点如果讲透，面试官会更容易相信你不是在“为了多一个名词而多一个组件”，而是在主动控制系统复杂度的来源。
\end{analysisBox}

\subsection{为什么要引入 Skill 体系，而不是只拼 Prompt}
\begin{center}
\small
\begin{tabularx}{\textwidth}{>{\bfseries}lXXX}
\toprule
方案 & 做法 & 优点 & 缺点 \\
\midrule
方案 A & 直接把一段 Markdown 或 system prompt 拼到 Planner 前面 & 实现快、改动小 & 只影响文案，不是体系对象，无法约束工具和 session \\
方案 B & 把 skill 解析成可注册对象，进入 session、planner、agent、tool allowlist & 能真正纳入体系，可解释、可审计、可热更新 & 需要额外的 parser、registry、CRUD 和状态同步 \\
\bottomrule
\end{tabularx}
\end{center}

\begin{qaBox}{最终选择}
如果目标只是“给模型多一点领域提示”，方案 A 就够了；但如果目标是“把领域能力做成系统级扩展点”，就必须选方案 B。当前项目已经按方案 B 落地：skill 以 Markdown + YAML frontmatter 形式存储，启动时加载进 runtime registry，session 会自动匹配并允许手动启用/禁用，最终同时影响 Planner、后续 Agent 提示和 tool allowlist。
\end{qaBox}

\begin{analysisBox}{深入分析}
这道题真正考的不是“你会不会做配置文件”，而是“你有没有把扩展能力做成一等系统对象”。如果只是把 skill 文件读出来，然后原样拼到 system prompt 里，它本质上还是一段文本，不是一个可治理对象。这样做最大的三个问题是：第一，session 不知道自己到底启用了哪些领域能力；第二，工具层不知道哪些 skill 允许哪些工具；第三，前端和审计层无法解释某次执行为什么走了这条策略。

至少有两个可行方案。方案 A 是 prompt-only，把 skill 当成额外文案。这种方式的优点是速度快、侵入性低，但它不会改变系统结构，最终也无法形成真正的 CRUD、热更新和 session 级开关。方案 B 是 registry 模式：skill 有稳定元数据、触发规则、允许工具、agent hints 和正文内容，运行时先做匹配，再把结果下沉到 session state。这样 Planner 可以按 skill 改写研究策略，Analyst / Reflection / Report 可以按 skill 加领域提示，工具层也能按 skill 收缩能力边界。这里必须选方案 B，因为“真正纳入体系”的定义，不是能被读出来，而是能影响调度、提示和权限。

为什么 skill 文件最终选择 Markdown + YAML frontmatter？因为这里同时有两类诉求：一类是机器可校验的稳定字段，例如 name、slug、trigger\_patterns、allowed\_tools；另一类是人类可编辑的正文说明，例如 overview、prompt、constraints。纯 JSON / YAML 对非工程用户不够友好，纯 Markdown 分节又太脆弱，frontmatter 则刚好把“结构化元数据”和“可读正文”分开。面试里如果你把这个点讲明白，面试官会知道你不是只在堆配置，而是在做一个面向运行时和面向人类编辑都可接受的能力层抽象。
\end{analysisBox}

\subsection{如果 Skill 规模很大，怎么避免启动慢、匹配慢和运行中断线}
\begin{qaBox}{结论}
海量 skill 的最优解，不是把所有 skill 全量读进内存，而是做成一套 \term{meta 常驻、content 延迟、匹配分层、session 固化、版本化回收} 的链路。这样既能处理冷启动和检索精度，也能保证 skill 更新或卸载后，运行中的 session 不断线。
\end{qaBox}

\begin{center}
\small
\begin{tabularx}{\textwidth}{>{\bfseries}lXXX}
\toprule
方案 & 做法 & 优点 & 缺点 \\
\midrule
方案 A & 启动时全量加载所有 skill 正文并全量扫描匹配 & 实现简单 & 冷启动慢、内存重、热更新抖、难扩展 \\
方案 B & meta/content 分层、候选渐进式匹配、正文按需加载、session snapshot 固化 & 冷启动和精度可平衡，运行时稳定 & 实现复杂度更高 \\
方案 C & 外部检索服务化，registry 只做路由 & 超大规模扩展性最好 & 过早分布式，运维成本高 \\
\bottomrule
\end{tabularx}
\end{center}

\begin{summaryBox}{最终选择}
当前最合理的是方案 B。因为它不要求一开始就把 skill 服务化，但已经能解决三个真实问题：第一，启动时不再依赖全量正文加载；第二，session 匹配时不再对所有 skill 做重规则扫描；第三，skill 更新和卸载后，新老 session 可以稳定分流，而不会把正在运行的任务打断。
\end{summaryBox}

\begin{analysisBox}{深入分析}
这个问题表面是在问“怎么做大规模加载”，但真正有三个耦合点。第一个是检索精度。如果你为了快，粗筛得太狠，就会漏掉真正该命中的 skill；如果你为了准，对所有 skill 做正则、规则和正文语义扫描，又会把延迟打爆。第二个是冷启动。如果启动阶段必须把所有 skill 的全文、正则和提示都预热进内存，实例扩容和服务重启会非常慢。第三个是状态连续性。如果运行中的 session 还依赖全局 registry 的当前态，那么 skill 一旦被更新、禁用或删除，就可能让同一个 session 在执行中前后行为漂移。

所以最优解必须同时回答这三件事。第一，在存储上做 meta/content 分离。meta 层只保留匹配和控制所需的轻字段，例如 slug、tenant、enabled、priority、keywords、allowed tools、agent hint 摘要；content 层保存完整 Markdown、解析后的 prompt 片段和较大的说明文本。启动时只加载 meta，不加载正文。第二，在匹配上做渐进式披露。先按 tenant、project、scope、enabled 做硬过滤，再用关键词、tag、domain 或全文索引做粗筛，把候选集从上万压到几十；然后只对候选集做正则、复杂 trigger rule 和更精细的策略合并；最后只为 top-N 命中的 skill 加载正文。这里“渐进式披露”不是前端折叠，而是计算成本和信息量逐级展开。第三，在运行态上做 session 固化。session 创建时生成 \term{resolved skill snapshot}，把 skill\_id、version、matched reason、effective tools、prompt fragments 都固化到会话里。之后 Planner、Analyst、Report 和 tool allowlist 全都只读这个 snapshot，而不是反复回查全局 registry。

这套设计为什么能解决 skill 卸载后不断线？关键在于版本固定和延迟回收。skill 更新不是原地覆盖，而是新版本；skill 卸载不是立刻把正文硬删，而是让它不再参与新 session 匹配。正在运行的 session 因为已经绑定了特定 version 的 snapshot，所以会继续沿用原版本执行。只有当没有 active session 再引用这个版本时，系统才真正回收正文内容。这就把“新流量不可见”和“旧流量不断线”分开了。

再往深一点说，海量 skill 的最大陷阱其实不是加载，而是 prompt 膨胀。即使你匹配很快，如果最后把二十几个 skill 的长正文全拼进 Planner 和 Report 的 prompt，LLM 侧还是会失控。所以真正成熟的系统还会再加一层约束：只让 top-N skill 进入执行，或把同类 skill 合并成摘要，并给 skill 注入本身设置 token budget 上限。面试里如果能把“匹配快”和“提示可控”一起讲出来，可信度会高很多。
\end{analysisBox}

\subsection{LLM 选型：为什么强调结构化输出和工具调用}
\begin{summaryBox}{选型标准}
这个系统里的模型不是单纯负责写文案，它要做 DAG 生成、JSON 结构化输出、事实校验、引用覆盖率分析和报告生成，所以第一优先级不是“最会聊天”，而是“最稳定地产生结构化结果”。
\end{summaryBox}

\begin{center}
\small
\begin{tabularx}{\textwidth}{>{\bfseries}lXXX}
\toprule
维度 & 倾向项 & 原因 & 对本项目的影响 \\
\midrule
工具调用稳定性 & 高优先级 & Planner / Reflection 都依赖结构化结果 & 降低解析失败和 fallback 频率 \\
中文能力 & 高优先级 & 用户查询、网页、报告以中文为主 & 减少跨语言语义损耗 \\
成本 & 高优先级 & 深研究任务 token 消耗天然高 & 决定是否能高频运行 \\
通用推理上限 & 中优先级 & 重要，但不是唯一约束 & 可由反思、检索和多轮流程补强 \\
\bottomrule
\end{tabularx}
\end{center}

\begin{qaBox}{面试时的稳妥表述}
我对模型选型的原则不是盲目追逐最强模型，而是优先看三个维度：结构化输出是否稳定、工具调用是否稳定、中文与成本是否平衡。因为在研究型 Agent 里，模型失败往往不是“文风不够好”，而是 JSON 不可解析、校验结果不稳定、成本高到没法长期运行。
\end{qaBox}

\subsection{浏览器自动化：为什么选择 Playwright}
\begin{center}
\small
\begin{tabularx}{\textwidth}{>{\bfseries}lXXXX}
\toprule
维度 & Playwright & Puppeteer & Selenium & 结论 \\
\midrule
异步模型 & 原生 async & 原生 async & 传统同步为主 & Playwright 更适合 Python async 服务 \\
自动等待 & 强 & 中 & 弱 & 减少页面交互竞态 \\
多浏览器支持 & Chromium/FF/WebKit & 偏 Chromium & 强 & Playwright 更平衡 \\
工程稳定性 & 高 & 中 & 中 & 对研究型抓取更合适 \\
\bottomrule
\end{tabularx}
\end{center}

\begin{qaBox}{核心回答}
我选 Playwright 的原因不是“它比较新”，而是它更贴合 Browser Use 场景。研究型浏览器节点经常面对动态加载、延迟渲染、滚动加载和反检测问题，Playwright 的自动等待、页面生命周期控制和无头执行能力更容易形成稳定工程实现。
\end{qaBox}

\subsection{为什么要实现三层浏览器提取策略}
\begin{qaBox}{设计理由}
不是所有页面都值得深度提取。浏览器节点如果默认把整页内容吃满，会直接推高 token 消耗、拉长时延、放大噪音。所以设计成 \term{snippet / skim / deep} 三层提取，本质上是把内容质量和成本一起治理。
\end{qaBox}

\begin{itemize}[leftmargin=2em]
  \item \term{snippet}：优先抓 meta 和少量核心文本，适合快速判断页面是否有价值。
  \item \term{skim}：抓主要段落，适合普通文章和资讯。
  \item \term{deep}：抓更多标题、段落、列表、结构化片段，适合权威报告和高价值页面。
\end{itemize}

\subsection{RAG 与向量库：为什么选择 pgvector}
\begin{center}
\small
\begin{tabularx}{\textwidth}{>{\bfseries}lXXXX}
\toprule
维度 & pgvector & Qdrant & Chroma & Pinecone \\
\midrule
部署复杂度 & 低 & 中 & 低 & SaaS \\
事务一致性 & 强 & 弱 & 弱 & 弱 \\
和业务表 JOIN & 方便 & 较弱 & 较弱 & 不方便 \\
全文搜索协同 & 强 & 需外部系统 & 弱 & 需外部系统 \\
运维成本 & 低 & 中 & 低 & 高 \\
\bottomrule
\end{tabularx}
\end{center}

\begin{qaBox}{最终选择}
我选 pgvector，不是因为它在纯向量召回上一定最强，而是因为它对这个项目的整体工程复杂度最低。文档、元数据、citation、research session、全文索引和向量索引都在同一套 PostgreSQL 里，事务一致性和数据联查能力非常强。
\end{qaBox}

\begin{analysisBox}{深入分析}
这里的关键不是“向量数据库哪个好”，而是“这个项目最需要优化的到底是哪一层”。如果你的系统是一个高吞吐、超大规模、纯向量召回优先的检索系统，那专用向量库可能更合理。但当前这个项目不是以 ANN 性能极限为第一目标，它更像一个研究工作流系统。研究工作流系统最常见的痛点不是“向量召回每秒少了几百次”，而是“文档、会话、引用、审计、预算、审批这些状态分散在多套存储里，最后查一次完整证据链要跨很多系统”。

pgvector 的优势正体现在这里。第一，它允许文档正文、结构化元数据、全文索引、向量索引、session 记录共库，从而把一致性问题压到最低。第二，很多研究场景并不是只靠向量召回，往往还要混合关键词搜索、过滤 group/source、回查 citation 元数据，PostgreSQL 在这些组合查询上非常顺手。第三，运维心智负担更低。对于一个还在快速演进治理闭环的系统来说，减少基础设施种类往往比追求单点最优性能更重要。

这题如果只回答“因为部署方便”会太浅。更好的讲法是：我不是在做一个独立的向量检索产品，我是在做一个带审计、带引用、带任务状态的研究系统，所以我优先优化的是系统整体复杂度，而不是某一个检索子模块的局部极值。
\end{analysisBox}

\subsection{为什么内部检索不是“查一下就回答”}
\begin{qaBox}{核心区分}
内部 RAG 在这个系统里不是回答器，而是证据源。它负责给出可引用片段、提供内部知识背景、帮助决定是否还要联网搜索，但最终回答仍然要经过 Analyst、Reflection 和 Report 链路。
\end{qaBox}

\subsection{生产追问：QA Agent 和 Deep Search Agent 的 SLO 为什么不同}
\begin{qaBox}{工程结论}
QA 业务型 Agent 的 SLO 是“确定边界内稳定、低时延、可控地答对”；Deep Search Agent 的 SLO 是“开放问题下覆盖充分、证据可追溯、冲突可解释”。这不是产品名差异，而是状态机、评估集、预算和失败语义都不同。
\end{qaBox}

\begin{analysisBox}{放回 DeepIntel 怎么讲}
这题如果脱离项目讲，很容易变成抽象概念题；放回 DeepIntel 会清楚很多。DeepIntel 更接近 Deep Search Agent，而不是单轮 QA Bot。它不是拿一个 query 直接查知识库然后输出，而是先由 Planner 拆成 DAG，再让 Search / Browser / RAG 分别去不同证据源取材，最后由 Analyst 做证据综合、Reflection 做质量校验，必要时还会 replan。也就是说，它的核心不是“快答”，而是“可验证地研究”。

至少有两种做法。方案 A 是做 QA 业务型 Agent：固定意图分类、固定工具路由、固定回答模板，优点是时延稳、维护成本低、SLA 好做，适合客服问答、内部知识问答、流程查询；缺点是开放问题时容易覆盖不足。方案 B 是做 Deep Search Agent：允许多跳检索、网页深读、多源比对和反思重规划，优点是复杂问题上更有上限，缺点是成本、时延和治理复杂度都更高。DeepIntel 选的是方案 B，因为它要解决的是“研究型任务”，而不是“FAQ 问答”。

面试时最稳的落点是：不要把两类 Agent 讲成谁替代谁，而要讲成优化目标不同。QA 业务型 Agent 更像高频、确定边界、强 SLA 的生产答题器；Deep Search Agent 更像低频、开放问题、证据优先的研究员。DeepIntel 当前主骨架明显站在后者这一侧。
\end{analysisBox}

\section{第三轮拷打：系统到底怎么跑，状态怎么流}

\begin{summaryBox}{这一轮面试官在试什么}
这一轮会从“你为什么这样设计”切到“这个系统运行时到底发生了什么”。如果你只能讲模块名，讲不清状态怎么变化、节点怎么路由、失败怎么传播，面试官会判断你更像读过代码，而不是做过系统。
\end{summaryBox}

\begin{qaBox}{这一轮最常见的追问}
\begin{itemize}[leftmargin=2em]
  \item 一个任务从 API 创建到报告产出，中间经过哪些确定的节点。
  \item Planner 产出的 DAG 和 LangGraph 本身是什么关系。
  \item public status 和 runtime status 为什么要分开，不分会出什么问题。
  \item 工具节点失败以后，状态怎么传给聚合器，聚合器怎么决定下一步。
\end{itemize}
\end{qaBox}

\subsection{从 API 到报告的完整生命周期}
\begin{center}
\input{figures/lifecycle}
\end{center}

\begin{summaryBox}{现状说明}
当前生命周期是：\term{create\_research} 建 session，状态先进入 \term{pending\_confirmation} 或 \term{running}，随后后台运行 LangGraph：\term{planner -> dag\_executor -> search/browser/rag -> dag\_aggregator -> analyst -> reflection -> replan/report -> END}。这条主链路已经打通，是系统最核心的业务骨架。
\end{summaryBox}

\subsection{生产追问：这个 Agent runtime 到底是 autonomous loop 还是显式 DAG}
\begin{qaBox}{工程结论}
DeepIntel 不是多 Agent 自由聊天，也不是单体 autonomous loop，而是 Planner 产显式 DAG，LangGraph 按依赖驱动 Search / Browser / RAG / MCP 节点，Harness 和 policy guardrail 在图外提供恢复、预算、审批和审计控制面。
\end{qaBox}

\begin{analysisBox}{按运行时语义拆开讲}
第一步是规划。用户问题进来后，不是立刻搜索，而是先由 Planner 把任务拆成一个显式 DAG，节点类型主要包括 search、browser、rag，另外还预留了 mcp。这里的重点不是“有个 JSON 计划”，而是依赖关系、并行性和节点类型都被结构化了，后续执行器才能知道哪些节点能一起跑、哪些必须等前置结果。

第二步是取证。LangGraph 里的 \term{dag\_executor} 会按可执行批次把节点分发给对应 Agent。Search 负责广撒网找入口，Browser 负责深读页面，RAG 负责查内部资料，而且项目里还显式做了 internal-first 约束：Web 节点默认要在内部 RAG 之后跑。也就是说，这不是“搜索一下再回答”，而是受图约束的并行取证。

第三步是综合。工具节点返回的不是最终答案，而是 evidence、tool history 和 node outcome。等 DAG 跑完后，Analyst 才把多源证据转成结构化分析；它会区分事实和推断、识别冲突、输出带 citation 的分析文本。第四步是校验与收口。Reflection 会把答案质量转成 \term{needs\_revision}、\term{overall\_confidence} 这类信号；如果校验不过且修订次数没到上限，就 replan 再走一轮，否则进入 Report。

方案对比也可以顺手讲一下。方案 A 是让多个 Agent 直接对话协作，优点是灵活，缺点是状态、依赖和恢复都很难治理。方案 B 是像 DeepIntel 这样，用 LangGraph + DAG 做显式编排，优点是执行顺序、预算、审批、审计和恢复都有稳定插点，缺点是前期建模更重。这个项目选方案 B，本质上是为了工程可控性。
\end{analysisBox}

\subsection{公开状态与内部运行状态为什么要分开}
\begin{qaBox}{核心设计}
对外 API 的状态要稳定，对内运行状态要足够细。公开层只需要 \term{pending / running / completed / failed} 这种粗粒度语义，但治理层必须区分 \term{awaiting\_approval / retryable\_failed / terminal\_failed / completed}，否则失败恢复、审批暂停和预算熔断都没法精确表达。
\end{qaBox}

\subsection{ResearchState 为什么是这个系统的核心}
\begin{lstlisting}[language=Python,caption={ResearchState 中最关键的状态维度}]
task_id / user_query / status / session
dag / current_executing_nodes / completed_nodes
tool_histories / node_outcomes / collected_evidence
verification / revision_needed / revision_count
runtime_status / budget_state / pending_approvals
agent_trace / guardrail_trace / errors
\end{lstlisting}

\begin{summaryBox}{设计原则}
ResearchState 不是简单的“共享上下文容器”，而是整个工作流的单一运行时状态机载体。凡是会影响路由、恢复、预算、审批和审计的字段，都应该显式出现在 state 里，而不是藏在零散的局部变量中。
\end{summaryBox}

\begin{analysisBox}{深入分析}
ResearchState 这一题本质上是在回答“系统的 source of execution truth 在哪里”。如果一个工作流系统没有统一的运行时状态载体，通常会出现三类问题。第一类问题是节点之间靠隐式约定通信，比如某个节点默认下游一定会去读某个临时变量或者某个 trace 片段，这会让路由和恢复逻辑变得脆弱。第二类问题是状态不完整，表面上流程能跑，但到了出错、回放、审计的时候才发现缺关键字段。第三类问题是治理字段没有一等公民地位，预算、审批、失败分类、runtime status 这类字段被塞进 \term{review\_status} 或普通文本里，后面很难做严肃的控制和恢复。

把这些状态提升到 ResearchState 层有一个很重要的工程收益：你可以把“路由条件”“恢复条件”“审计快照”“预算熔断”都建模为状态上的显式变换，而不是 if-else 里的旁路副作用。这样做的代价当然也有，就是 state 会逐步变大、会需要整理 merge 语义、会需要对哪些字段累加、哪些字段覆盖有严格约束。但对于一个长期任务系统来说，这个代价是值得付的，因为它换来的是可解释性和可恢复性。
\end{analysisBox}

\subsection{Planner 如何生成 DAG，为什么还需要 fallback}
\begin{qaBox}{面试回答}
Planner 的职责不是直接给最终答案，而是把用户问题拆成一张可执行研究图。它会让 LLM 生成 \term{DAGDefinition}，里面包含节点类型、查询内容、依赖关系和并行性。但因为 LLM 可能输出坏 JSON、空结果或不合理节点，所以必须准备 fallback DAG，保证流程至少能退化成“搜索 + 浏览 + 内部检索 + 分析”这条保底路径。
\end{qaBox}

\begin{analysisBox}{深入分析}
Planner 这一题考的不是“会不会让模型吐 JSON”，而是“你有没有把 LLM 放在一个可失败的位置上”。如果把 Planner 产物当绝对真理，系统会在两个地方变脆。第一，结构脆弱。模型可能输出坏 JSON、漏依赖、生成不存在的节点类型，或者把本该并行的步骤串起来。第二，策略脆弱。即便 JSON 结构合法，它也可能给出一个业务上很差的计划，比如过度依赖外网、把低价值网页抓取得过深、或者在证据不足时就安排报告节点。

所以 fallback 的意义不是给 Planner 擦屁股，而是承认 Planner 是智能组件，不是可信编译器。工程上至少有两个方案。方案 A 是 Planner 失败就让整个任务失败，优点是语义简单，缺点是把整个系统可用性绑死在一次模型输出上。方案 B 是 Planner 尽量生成高质量 DAG，但执行层永远保留一条保底研究路径。这里显然要选方案 B，因为研究系统的首要目标是“任务可继续”，而不是“计划一定漂亮”。

更深一层看，fallback 还承担了平台语义稳定器的角色。它把最低可执行能力固定成一套可预测的研究骨架，使得即便未来换模型、换 prompt、换节点模板，系统也不会因为 Planner 一次抖动就彻底失能。面试时如果你能把 fallback 讲成“系统韧性设计”，而不是“异常处理补丁”，说服力会强很多。
\end{analysisBox}

\begin{qaBox}{如果面试官追问“未知问题下灵活性会不会为零”}
不会。这里的灵活性不是靠预先写死所有备选方案，而是靠三层机制叠加出来的。第一层是 Planner 动态生成 DAG，而不是固定流程。第二层是 Planner 失败时有 fallback DAG，保证任务不会因为一次模型输出坏掉就整体失能。第三层是 Reflection 和 replan 形成闭环，如果第一次计划覆盖不够或者执行中暴露出失败模式，系统会带着失败记忆重新规划，而不是机械重跑旧图。
\end{qaBox}

\begin{qaBox}{原因}
很多人会把“有 Planner”误解成“已经有动态规划能力”，这不够准确。真正的动态规划能力至少要回答三个问题。第一，第一次计划失败怎么办；第二，执行到一半发现路走错了怎么办；第三，同样的错误再次出现时，系统会不会还按原路硬撞。如果这些问题答不上来，未知问题下的灵活性确实会很低。
\end{qaBox}

\begin{qaBox}{选型}
这里至少有两个方案。方案 A 是固定流程加局部兜底，例如 Search 失败就换 Browser，或者 Planner 失败就直接终止。优点是实现简单，缺点是备选路径是预写死的，面对未知问题很快失效。方案 B 是 Planner 动态生成 DAG，执行层保留 fallback DAG，质量层再通过 Reflection 和 replan 形成带记忆的闭环。这里必须选方案 B，因为研究型任务的关键不是提前枚举所有备选方案，而是让系统在运行中根据失败和证据情况重组路径。
\end{qaBox}

\begin{qaBox}{如何做}
当前这套实现里，第一次计划失败会退化到 fallback DAG，保证最小可执行能力不丢。执行到一半如果 Reflection 认为证据不足、引用覆盖不够或者整体置信度低，系统不会只把失败节点重置成 pending，而是会进入 replan，重新调用 Planner 生成一张新 DAG。更重要的是，replan 不是在真空中发生的，它会把本会话已经出现过的失败模式整理成 failure memory 注入新的规划提示，明确告诉 Planner 哪些工具模式、查询模式或错误类别已经踩过坑，优先换工具、缩小查询范围或改变取证路径。
\end{qaBox}

\begin{qaBox}{效果}
这样处理以后，系统的灵活性不再依赖预先写死多少分支，而是依赖三种可组合能力：动态拆图、保底退化、带记忆重规划。它解决的不是“Planner 永远正确”，而是“Planner 不正确时任务还能继续，而且第二次规划会比第一次更少重复犯错”。
\end{qaBox}

\subsection{DAG 的正确性如何保证}
\begin{itemize}[leftmargin=2em]
  \item 数据结构层：Planner 输出被解析成 \term{PlanNode / PlanEdge / DAGDefinition}。
  \item 执行层：\term{get\_executable\_order()} 用拓扑思路分批求可执行节点。
  \item 容错层：非法节点会被过滤，Planner 失败时回退到 fallback DAG。
  \item 策略层：系统还会强制内部 RAG 先于 Web 节点执行，形成 internal-first 检索策略。
\end{itemize}

\begin{analysisBox}{深入分析}
DAG 正确性不能只理解成“有没有环”。在这个系统里，它至少有三层含义。第一层是结构正确，即节点类型合法、依赖存在、拓扑可排序。第二层是执行正确，即调度器不会把失败节点当成完成，也不会在依赖未满足时提前放行下游。第三层是策略正确，即图虽然形式上可跑，但是否符合平台要求，例如 internal-first、预算优先、审批前不得执行高风险写操作。

这也是为什么单靠 Planner 输出校验不够。很多人会说“把 JSON schema 校验做好就行”，这只能保证结构层。真正难的是执行层和策略层，因为这些正确性是运行时才能暴露的。比如某个 Browser 节点虽然结构合法，但它连续超时后是否应该进入 retryable failed，还是被预算 breaker 直接跳过；某个 MCP 写工具虽然在图里合法，但如果缺审批，运行时就必须被挂起。这些都不是静态 schema 能解决的问题。

所以这里的正确性设计本质上是“静态约束 + 动态约束”的组合。静态部分让 DAG 至少能被解释，动态部分保证图在真实世界里按治理要求运行。面试里如果你能明确这两层，说明你理解的是工程正确性，而不是算法课里的 DAG 定义。
\end{analysisBox}

\subsection{为什么要先 internal RAG，再决定要不要联网}
\begin{qaBox}{设计理由}
研究型系统不应该把互联网搜索当作唯一真相源。内部知识库通常更贴近企业语境，也更可控，所以系统会优先执行内部 RAG，再根据命中情况和策略判断是否继续走 Web。这样既减少无意义联网，也能降低把内部问题完全外部化的风险。
\end{qaBox}

\subsection{工具节点为什么必须返回显式结果契约}
\begin{lstlisting}[language=Python,caption={NodeOutcome 契约}]
NodeOutcome = TypedDict("NodeOutcome", {
    "node_id": str,
    "tool_call_id": str | None,
    "tool_name": str,
    "status": Literal[
        "success",
        "retryable_error",
        "terminal_error",
        "awaiting_approval",
        "skipped",
    ],
    "error_category": str | None,
    "error_message": str | None,
    "retry_count": int,
    "tokens_used": int,
    "cost_usd": float,
    "result_count": int,
    "approval_request_id": str | None,
})
\end{lstlisting}

\begin{qaBox}{为什么不能只靠 tool\_histories}
\term{tool\_histories} 更适合审计，不适合做唯一状态源。因为审计记录的目标是“发生了什么”，而路由的目标是“下一步怎么办”。如果失败语义只藏在 tool history 里，聚合器很容易在语义不明确时把节点误标成完成。显式的 \term{NodeOutcome} 正是为了解决这个问题。
\end{qaBox}

\begin{analysisBox}{深入分析}
这道题是整个治理改造里最有代表性的一个点，因为它直接说明“记录”和“控制”不能混为一谈。\term{tool\_histories} 天然是面向审计的结构，它关注的是调用了什么工具、传了什么参数、耗时多久、结果摘要是什么。它当然很重要，但它不是为路由设计的。路由真正需要的，是一组最小而明确的控制语义：这次调用是成功、可重试失败、终止失败、等待审批，还是被跳过；如果失败，失败类别是什么；如果等待审批，审批请求 ID 是什么。

如果没有 NodeOutcome，聚合器通常只能从一堆宽泛的历史记录中“推断”当前节点是不是算完成。这种推断最危险的地方在于，一旦工具层记录不完整或者语义不稳定，执行层就会把错误状态误判成完成状态。历史上那个“工具失败但节点被标 done”的 bug，本质上就是因为控制语义没有一等建模。

所以 NodeOutcome 的价值不在于字段好看，而在于它把执行层真正需要的控制信息抽了出来。你可以把它理解成工具层和调度层之间的一份契约：工具层负责把业务结果和失败语义翻译成一组最小控制信号，调度层再用这组信号决定下一步如何路由。只有这样，审计和控制才不会互相污染。
\end{analysisBox}

\subsection{怎样避免单点故障把整条链路拖死}
\begin{center}
\small
\begin{tabularx}{\textwidth}{>{\bfseries}lXXX}
\toprule
方案 & 做法 & 优点 & 缺点 \\
\midrule
方案 A & 任意节点失败就整任务失败 & 语义简单 & 单点故障风险高，未知问题可用性差 \\
方案 B & 失败显式分类，临时故障指数退避重试，永久故障终止或重规划 & 可恢复性强，治理语义清晰 & 需要更多状态和策略维护 \\
\bottomrule
\end{tabularx}
\end{center}

\begin{qaBox}{原因}
如果系统不做错误分类，最典型的后果就是两头都错。第一类错法是把瞬时抖动当永久错误，明明一次 retry 就能恢复，却直接把整条链路打死。第二类错法是把永久参数错误当临时波动，结果系统拿着错误输入不断重试，把预算和 wall clock 一起烧掉。更进一步，如果可重试失败每次都立刻再打一次，上游正处在抖动窗口时，系统只会更快撞上限流和预算阈值。
\end{qaBox}

\begin{qaBox}{选型}
这里至少有三个可行方案。方案 A 是所有失败都统一成一个布尔值，只决定成功或失败，实现最简单，但治理语义几乎为零。方案 B 是显式区分 \term{retryable\_error} 和 \term{terminal\_error}，但失败后立即有限重试，优点是比方案 A 强很多，缺点是对上游抖动不够友好。方案 C 是显式失败分类，再把可恢复错误交给调度层做指数退避重试，同时把终止性错误直接停掉。这里要选方案 C，因为它既能区分错误性质，又能避免在抖动窗口里无意义地密集重试。
\end{qaBox}

\begin{qaBox}{如何做}
当前实现里，工具节点不会再把所有失败都压成一个布尔值，而是显式区分 \term{retryable\_error}、\term{terminal\_error} 和 \term{awaiting\_approval}。timeout、429、上游 5xx、网络抖动、临时 DB 故障会被归到可恢复类别；schema 校验失败、权限问题、审批拒绝、MCP server 不可信这类会被归到终止类别。对于前一类错误，系统不会立刻机械重试，而是根据错误类别给一个基础退避时间，再随着 \term{retry\_count} 指数增长，计算出 \term{next\_retry\_at} 写进 failure memory，并受 \term{max\_retries\_per\_tool} 约束。调度层只有在退避窗口到期后，才会把节点重新放回执行队列。对于后一类错误，节点会直接标成 terminal，不再重复消耗预算，同时失败模式会写入 failure memory，为后续 replan 提供负反馈。
\end{qaBox}

\begin{qaBox}{效果}
这样做以后，系统既不会因为一次网络抖动就把整条任务打死，也不会因为一个永久参数错误而在错误输入上疯狂重试。可恢复错误会在调度层自动退避并等待更合适的重试窗口，终止性错误会被尽快收口，避免继续烧 token、tool call 和 wall clock。
\end{qaBox}

\subsection{系统会不会从历史失败里吸取教训}
\begin{qaBox}{原因}
如果系统每次失败都把历史清空，就会重复踩同一种坑。对研究型任务来说，真正浪费预算的往往不是单次失败，而是同一个工具模式、同一个查询模式、同一种错误类别在一个会话里反复出现。这个问题如果不处理，重试会退化成“重复犯错”，replan 也会退化成“换个 prompt 再犯一遍”。
\end{qaBox}

\begin{qaBox}{选型}
这里至少有两个方案。方案 A 是完全不记失败历史，每次重试和 replan 都只看当前节点状态，优点是实现简单，缺点是系统没有记忆。方案 B 是先做会话级 failure memory，把本次任务里的失败模式变成控制面的一等状态，再用于重试调度和 replan 提示。这里选方案 B，而不是一上来做全局跨任务学习，是因为会话级记忆最贴近当前控制面，也最容易验证效果和控制风险。
\end{qaBox}

\begin{qaBox}{如何做}
当前已经支持会话级 failure memory，而不是全局跨会话学习。系统会总结本次任务里哪些工具模式、查询模式和错误类别已经失败过，并把这些信息记录成结构化状态。它至少做三件事。第一，重复失败阻断。同一个工具类型加同一个查询，如果已经达到最大重试次数，或者已经被判成终止性失败，系统不会继续把它当成普通 pending 节点调度。第二，退避调度。对于可恢复错误，系统会把下一次可重试时间写进 failure memory，按指数退避窗口再入队。第三，重规划提示。系统会把这些失败模式整理成简短的 planning hints 注入下一次 replan，明确告诉 Planner：哪些模式已经失败过、哪些已经 repeat blocked、优先换工具还是缩窄查询。这里也要守住边界：当前退避是图内调度层的短等待与再入队，不是 Harness 级定时唤醒系统。
\end{qaBox}

\begin{qaBox}{效果}
failure memory 的效果不是把系统吹成“会自我进化的大模型”，而是先把失败经验转化成可执行控制。它能减少会话内重复犯错，让自动重试更克制，让 replan 更有方向感。面试里最稳的说法是：我们已经实现了会话级失败记忆，能在本次任务里做到重复失败阻断、指数退避调度和带记忆重规划；但还没有把它做成跨任务长期学习系统。
\end{qaBox}

\subsection{Search / Browser / RAG 三类节点如何分工}
\begin{itemize}[leftmargin=2em]
  \item Search：适合快速事实、数据、新闻入口、网址发现。
  \item Browser：适合深读文章、结构化页面提取、动态网页。
  \item RAG：适合内部资料、历史文档、背景知识和企业语义上下文。
\end{itemize}

\begin{summaryBox}{面试回答模板}
这三类节点不是冗余，而是故意分层。Search 负责广撒网，Browser 负责深挖高价值页面，RAG 负责内部可控证据。真正有价值的研究系统不是任何问题都扔给浏览器，而是先判断哪种证据源最值得花成本。
\end{summaryBox}

\section{第四轮拷打：关键机制是不是你真的做过}

\begin{summaryBox}{这一轮面试官在试什么}
这一轮会专门挑几个最容易暴露“只懂概念、不懂实现”的点来问，例如 Reflection 有什么用、RRF 为什么合理、Checkpointer 到底帮你解决了什么。这里需要把抽象概念落回具体执行语义。
\end{summaryBox}

\begin{qaBox}{这一轮最常见的追问}
\begin{itemize}[leftmargin=2em]
  \item Reflection 是不是只是让模型再看一遍自己的答案。
  \item Analyst 和 Report 为什么不能合并成一步。
  \item Checkpointer 到底恢复了什么，没恢复什么。
  \item SSE 只是前端体验，还是系统治理的一部分。
\end{itemize}
\end{qaBox}

\subsection{Analyst 为什么单独存在}
\begin{qaBox}{核心作用}
Analyst 的意义是把多源证据从“片段集合”转成“结构化分析”。Search、Browser、RAG 拿回来的是证据，Analyst 负责做事实归类、冲突识别、分析组织和初步结论形成。
\end{qaBox}

\begin{analysisBox}{深入分析}
这一题真正要回答的是，为什么不能让最终 Report 一步做完。答案是证据整合和结果表达虽然都要调用模型，但它们承担的是不同职责。Analyst 关注的是“证据之间到底是什么关系”，包括哪些事实互相印证、哪些来源互相冲突、哪些结论目前证据不足；Report 关注的是“如何把已经形成的分析组织成适合阅读的输出”。把两者合并，短期能省一个节点，长期会让分析过程变成黑箱。

分开设计至少有三个收益。第一，错误隔离。若报告写得不好，你可以判断是表达问题还是前面的分析就错了。第二，治理插点更清晰。Reflection 应该审的是分析质量和引用覆盖，而不是一篇已经高度修辞化的最终文案。第三，复用能力更强。未来如果要同时输出简版摘要、详细报告、内部 memo，Analyst 产出的结构化分析可以被多个 report renderer 复用。

所以 Analyst 不是为了“多一个 Agent”，而是为了把证据解释层从最终呈现层剥离出来。面试官如果继续追问“是不是过度设计”，你就可以直接回答：只有当系统目标是一次性回答时它才显得重；一旦目标是多源证据研究、复盘和质量控制，这个拆分就是必要的。
\end{analysisBox}

\subsection{Reflection 为什么不能省掉}
\begin{qaBox}{核心回答}
如果没有 Reflection，整个系统就会退化成“拿到一些证据后直接写报告”。Reflection 的职责是把回答质量结构化成 \term{overall\_confidence}、\term{citation\_coverage}、\term{needs\_revision} 等信号，再由图路由决定是否重规划。它不是装饰，而是把“能不能继续往前走”显式化。
\end{qaBox}

\begin{analysisBox}{深入分析}
Reflection 的关键价值，是把“质量判断”从人脑直觉变成系统信号。没有它，系统虽然也能输出报告，但整个链路没有一个地方正式回答“这份报告到底靠不靠谱、证据覆盖够不够、是否需要再搜一次”。很多 Demo 都止步在这一步，因为只要模型能写出看起来像样的报告，就误以为流程闭环了。实际上那只是生成闭环，不是质量闭环。

备选方案至少有两个。方案 A 是完全不做 Reflection，相信 Analyst 和 Report 已经足够。优点是链路短、成本低，缺点是系统没有内生的质量回路。方案 B 是加入一个专门的反思节点，把答案质量、引用覆盖率和修订建议结构化，再把它转成路由条件。这里必须选方案 B，因为研究系统的价值不只是“能写”，而是“在证据不足时知道自己该停下还是重来”。

当然，Reflection 也不是银弹。它本身仍依赖模型判断，所以它给出的不是绝对真相，而是一个可操作的质量信号层。这一点要诚实讲清楚。它不是替代人工 reviewer，而是在自动链路里尽量提前发现低置信或低覆盖的问题。这样表达既不会夸大能力，也能体现你理解它的工程位置。
\end{analysisBox}

\subsection{Reflection 和传统 reviewer 的区别}
\begin{itemize}[leftmargin=2em]
  \item Reviewer 更偏规则门禁或人工审查。
  \item Reflection 更偏模型辅助的自校验与质量评分。
  \item 在当前项目里，Reflection 和动作审批 reviewer 已经分开：Reflection 负责答案质量自校验，approval reviewer 负责高风险动作的暂停、审批和恢复入口。
\end{itemize}

\subsection{RRF 融合检索为什么合理}
\begin{lstlisting}[language=Python,caption={RRF 的直观公式}]
score(d) = sum(1 / (k + rank(d)))
\end{lstlisting}

\begin{summaryBox}{设计理由}
RRF 不是为了追求算法花哨，而是为了在多路检索结果之间取得稳定折中。向量相似度、关键词匹配、重排结果的分值尺度往往不一致，直接线性加权很脆弱；RRF 只依赖 rank，更稳、更简单，也更容易工程落地。
\end{summaryBox}

\begin{analysisBox}{深入分析}
RRF 合理，不是因为它在论文里很经典，而是因为它非常适合工程环境里的“异构召回”。在真实系统里，向量检索、关键词检索、规则检索、重排模型往往来自不同模块，它们的分值空间根本不可直接比较。你如果做线性加权，就必须不断调分值归一化、权重比例、场景偏置，这在数据分布变化时会很脆。

RRF 退一步，只相信名次，不相信原始分数。这样做牺牲了一些极致调优空间，但换来的是稳定性和可维护性。对于一个研究型系统来说，这通常是正确取舍，因为它需要一个在数据源变化、召回器变化时依然不至于失控的融合策略，而不是一个只在离线集上分数最好、上线后难维护的复杂公式。

如果面试官追问为什么不用学习排序或者更复杂的融合，你可以直接回答：可以做，但前提是你已经有足够稳定的数据闭环和标注反馈。当前阶段更值得优先的是稳健和可解释，而不是把排序问题过早做成重机器学习工程。这个回答会显得很务实。
\end{analysisBox}

\subsection{Checkpointer 的真实价值是什么}
\begin{qaBox}{核心答案}
\term{thread\_id=session\_id} 保证的是会话隔离，Checkpointer 提供的是状态恢复能力。它能让 LangGraph 在节点边界保存执行快照，使流程在中断后有机会从已知状态继续。但要注意，检查点本身不是完整治理方案，它不自动解决预算、审批、审计和 supervisor 恢复控制。
\end{qaBox}

\begin{analysisBox}{深入分析}
这一题最常见的误区是把 checkpointer 神化，好像一旦有了 checkpoint，恢复问题就全解决了。实际上 checkpointer 解决的是“图内状态快照”问题，不是“任务级治理闭环”问题。它很擅长保存某个 thread 在某个节点边界上的状态，让你不必每次都从头开始跑；但它不天然知道当前 worker 是否还持有租约，不知道预算是否已超，不知道某个动作是否还在等待人工审批，也不知道 supervisor 现在是否决定暂停整个任务。

所以 checkpointer 的正确定位是：它是恢复机制的重要组成部分，但不是完整的恢复控制面。你可以把它理解成恢复所需的“状态底座”，而 Harness 和数据库里的 runtime/approval/budget 状态则是在定义“到底能不能恢复、应该恢复到哪里、恢复后是否继续执行”。面试时如果你能把 checkpointer 的能力边界讲清楚，反而更显得你做过真实长任务系统，而不是把框架能力误当成系统能力。
\end{analysisBox}

\subsection{为什么 SSE 对这个系统重要}
\begin{qaBox}{核心理由}
研究任务天然是长任务，如果没有流式可观测性，用户只能看到一个很长的 loading。SSE 的价值不是“酷炫实时输出”，而是把 agent trace、tool 进度、状态变化、报告片段和错误事件透明推给前端，让任务变成可感知、可追踪、可中断观察的过程。
\end{qaBox}

\section{第五轮拷打：性能、成本和资源怎么扛}

\begin{summaryBox}{这一轮面试官在试什么}
很多候选人在这里会只讲“做了并发、做了缓存、做了优化”。这不够。研究型 Agent 的资源问题至少包含三类：时延、token/cost、数据库与工具开销。面试官想看的是你有没有把它们视为同一个系统问题。
\end{summaryBox}

\begin{qaBox}{这一轮最常见的追问}
\begin{itemize}[leftmargin=2em]
  \item 这个系统最主要的耗时在什么地方，是模型、检索还是网页抓取。
  \item 你怎么避免 Browser 节点把 token 直接吃爆。
  \item 预算治理是软提示还是硬熔断，熔断之后系统怎么收口。
  \item 为什么数据库仍然选 PostgreSQL，而不是拆成更多专用组件。
\end{itemize}
\end{qaBox}

\subsection{为什么并行执行能显著降低时延}
\begin{lstlisting}[language=Python,caption={LangGraph Send API 的作用}]
dag_executor -> [Send("search"), Send("browser"), Send("rag")]
\end{lstlisting}

\begin{qaBox}{面试回答}
如果 Search、Browser、RAG 串行执行，总时延接近三者之和；并行 fan-out 后，总时延更接近三者中的最大值。这类任务的瓶颈往往在 I/O，而不是 CPU，所以并行调度的收益很高。
\end{qaBox}

\begin{analysisBox}{深入分析}
并行执行的收益不能只用一句“总时延取最大值”带过，面试官真正想看的是你是否知道什么场景适合并行、什么场景不适合。这个系统里 Search、Browser、RAG 的共同特点是高度 I/O 密集，而且很多时候相互独立。Search 等网络返回，Browser 等页面加载，RAG 等数据库检索，这些等待时间如果串行堆起来，就是纯浪费。所以把它们 fan-out，本质上是在回收等待时间。

但并行不是越多越好。第一，外部依赖会有限流和配额，Browser 开太多页面可能直接把系统拖死。第二，并行会放大预算消耗，如果 Planner 过度拆图，虽然时延下降了，但 token、工具调用次数和外部 API 成本会同时上升。第三，结果聚合会更复杂，失败语义、超时处理和 partial result 策略都必须更严谨。

所以正确回答不是“并行一定更快”，而是“在 I/O 主导、节点相互独立、预算允许的前提下，并行能显著降低 wall clock；但必须配合批次调度、预算 breaker 和失败聚合语义，否则它只是把串行浪费变成并行失控”。这句话会让面试官感觉你真的踩过并发系统的坑。
\end{analysisBox}

\subsection{Token 优化策略}
\begin{itemize}[leftmargin=2em]
  \item 输出长度分 \term{short / medium / long} 三档，不只是影响报告 token，也同步影响搜索次数、RAG 返回量和 citation 上限。
  \item 浏览器内容用 \term{snippet / skim / deep} 控制页面读取粒度。
  \item RAG 侧采用候选召回 + 排序 + 截断，避免把低价值片段灌入上下文。
  \item 最终报告阶段再根据预算控制输出规模，避免前面省下来的 token 最后一次性烧光。
\end{itemize}

\begin{analysisBox}{深入分析}
Token 优化真正难的地方，不在于“少输出一点”，而在于你不能为了省 token 把证据链和答案质量一起省没了。研究型 Agent 和普通聊天模型的最大区别在于，很多 token 消耗不是出在最终回答，而是出在前面的证据获取和中间分析。网页抓取太深、RAG 返回太多低质量片段、Planner 把 DAG 拆得太碎、Reflection 重复审查太多次，都会在最终报告生成前就把预算吃掉。

所以这里的优化必须是系统性的。第一层是输入压缩，即在进入模型前先控制证据粒度，例如浏览器的 snippet/skim/deep、RAG 的候选截断、citation 上限。第二层是路径裁剪，即并不是所有任务都值得走完整 Web 扩展链路，internal-first 的策略就是一种成本治理。第三层才是输出控制，即最终报告的长度、引用数量、展开程度。把这三层分开讲，面试官会更容易相信你理解的是“成本结构”，而不是只有“把 max\_tokens 调小”这种表层操作。
\end{analysisBox}

\subsection{数据库优化怎么看}
\begin{itemize}[leftmargin=2em]
  \item 文档层：向量索引承担 ANN，全文索引承担关键词检索。
  \item 业务层：research session、citation、tool audit、budget state 都在同一库里，降低同步复杂度。
  \item 可观测层：指标与 trace 可以按 session 维度做统计和聚合。
\end{itemize}

\begin{analysisBox}{深入分析}
数据库选型这题不能只答“PostgreSQL 通用、成熟”。真正的判断是：当前系统的复杂度是否值得拆成多套专用存储。至少有两个方案。方案 A 是 PostgreSQL 一库承载业务数据、向量检索、审计和预算状态，优点是事务边界清晰、查询方便、运维成本低，缺点是向量能力和高吞吐日志能力不如专用系统。方案 B 是把向量库、OLTP 库、日志库拆开，优点是各组件可以针对性优化，缺点是同步、幂等、一致性和运维复杂度都会上升。

当前阶段选方案 A 是合理的，因为系统最主要的矛盾还不是极致吞吐，而是治理闭环和状态一致性。审批、预算、runtime status、tool audit 这些关键控制数据如果分散在多种存储里，恢复和审计会先变复杂。把它们尽量放在同一个事务型主库里，反而更利于把控制面做稳。

当然，这个选择有边界。随着文档规模、检索 QPS、审计写入量上升，未来完全可能拆出专用检索或时序/日志存储。但那应该建立在当前单库模型已经成为明确瓶颈之后，而不是一开始为了“架构高级感”就过度拆分。面试时把这一层演进逻辑说清楚，会比直接喊“微服务 + 多存储”成熟得多。
\end{analysisBox}

\section{第六轮拷打：你怎么证明它没胡来}

\begin{summaryBox}{这一轮面试官在试什么}
会做流程不等于会做系统。面试官在这一轮通常会问：你怎么知道它真的在工作，而不是表面有输出？你怎么定位问题？你怎么区分是模型错了、工具错了、还是治理规则拦住了？
\end{summaryBox}

\begin{qaBox}{这一轮最常见的追问}
\begin{itemize}[leftmargin=2em]
  \item 你会盯哪些指标，而不是泛泛地说“看日志”。
  \item trace 为什么要拆成 agent、guardrail、tool audit 三条线。
  \item 出现错误时，怎么判断是模型输出坏了还是外部工具失败了。
  \item 这个系统能不能做法证级复盘，复盘数据现在放在哪。
\end{itemize}
\end{qaBox}

\subsection{我会看哪些指标}
\begin{center}
\small
\begin{tabularx}{\textwidth}{>{\bfseries}lXX}
\toprule
维度 & 代表指标 & 作用 \\
\midrule
工作流层 & 成功率、P95 时长、重规划次数 & 看流程稳定性 \\
工具层 & tool success rate、平均时长、错误分类 & 看节点可靠性 \\
质量层 & 离线看 Hit@K、MRR、Faithfulness、Answer Relevancy；在线看 citation coverage、overall confidence、转人工率 & 分层定位检索问题还是生成问题 \\
成本层 & used tokens、used cost、tool calls、wall clock & 看预算消耗 \\
治理层 & approval count、budget stop count、terminal failures & 看可控性 \\
\bottomrule
\end{tabularx}
\end{center}

\begin{summaryBox}{面试里这题的标准说法}
RAG 评估不能只看最终答案好不好，而要拆成两层。第一层先单独看检索层，常用 Hit@K 和 MRR。Hit@K 回答“正确 chunk 找到没”，MRR 回答“找到得靠不靠前”。第二层再看生成层，常用 Faithfulness 和 Answer Relevancy，前者看有没有照着资料说，后者看有没有回答用户真正的问题。线上再用点击率、转人工率这类业务指标做最终验收。这样一拆，问题才能定位得准。
\end{summaryBox}

\begin{figure}[htbp]
  \centering
  \includegraphics[width=0.92\textwidth]{figures/rag-two-layer-framework.jpg}
  \caption{RAG 两层评估框架：离线先拆检索层与生成层，线上再看业务验收指标}
\end{figure}

\begin{analysisBox}{深入分析}
如果只盯最终答案质量，很容易把检索问题和生成问题混在一起。比如答案不好，可能是检索根本没把正确 chunk 召回来，也可能是 chunk 明明召回了，但排序太靠后，模型没好好用；还可能是上下文没问题，但模型自己跑题或者幻觉。把评估拆成检索层和生成层，诊断路径就会清晰很多。

检索层里，Hit@K 看的是“有没有找着”。它不关心正确 chunk 是第一名还是第五名，只关心在 Top-K 里有没有出现。所以 Hit@5 高，说明召回大体没丢。MRR 看的是“找着得有多快”。如果 Hit@5 很高，但 MRR 很低，通常表示正确内容虽然被召回了，但排位靠后，这时候优先怀疑 rerank、融合排序或者 chunk 粒度不合适。

生成层里，Faithfulness 和 Answer Relevancy 不能混为一谈。Faithfulness 问的是“说的是不是真的，有没有证据支撑”，Answer Relevancy 问的是“说的是不是用户想要的”。一个答案可以非常忠实于资料，但完全答非所问；也可以回答了问题，但中间夹杂很多编造内容。面试时把这两个维度分开讲，面试官通常会认为你是真的做过质量诊断，而不是只背了几个名词。

工程上我会把这套体系分成离线和在线两条线。离线用小规模标注集跑 Hit@K、MRR，以及基于 RAGAs 思路的 Faithfulness、Answer Relevancy、Context Recall、Context Precision，用来快速定位是哪一层出了问题。在线则继续看 citation coverage、overall confidence、点击率、转人工率，把它们作为最终业务验收。这样既能做研发调参，也能做真实效果闭环。
\end{analysisBox}

\begin{figure}[htbp]
  \centering
  \includegraphics[width=0.78\textwidth]{figures/hitk-vs-mrr.jpg}
  \caption{Hit@K vs MRR：Hit 看“有没有找到”，MRR 看“找到得有多靠前”}
\end{figure}

\begin{qaBox}{怎么解释 Hit@K 和 MRR 的差异}
很多人会把 Hit@K 和 MRR 讲成一回事，这不够准确。Hit@K 只回答“Top-K 里有没有命中正确 chunk”，所以它更像召回覆盖指标；MRR 则会惩罚排名靠后，正确内容排第一名得 1 分，排第五名只有 0.2 分，更适合判断 rerank 和融合排序有没有把关键内容顶到前面。面试时最好直接举一个例子：Hit@5 很高但 MRR 很低，说明“找到了，但排得不够前”。
\end{qaBox}

\begin{figure}[htbp]
  \centering
  \includegraphics[width=0.8\textwidth]{figures/faithfulness-vs-answer-relevancy.jpg}
  \caption{Faithfulness vs Answer Relevancy：忠实于资料不等于回答了用户问题}
\end{figure}

\begin{qaBox}{怎么解释 Faithfulness 和 Answer Relevancy}
Faithfulness 和 Answer Relevancy 不能混着说。Faithfulness 关注答案是否有证据出处，本质上是在看幻觉；Answer Relevancy 关注答案是不是在回答用户真正的问题，本质上是在看跑题。一个回答可以字字有据但完全答非所问，也可以看上去回答了问题但夹杂大量编造内容。把这两个维度拆开，生成层的问题才能定位得准。
\end{qaBox}

\begin{figure}[htbp]
  \centering
  \includegraphics[width=0.92\textwidth]{figures/rag-diagnosis-map.jpg}
  \caption{RAG 指标诊断与优化映射：不同指标低，对应不同病灶和优化方向}
\end{figure}

\begin{analysisBox}{指标低时怎么诊断}
这张图在面试里特别好用，因为它能把“会看指标”和“会做优化”连起来。如果 Context Recall 低，优先怀疑 embedding、chunking、多路召回不够；如果 Context Precision 低，优先怀疑 rerank 或 Top-K 截断不合适；如果 Faithfulness 低，优先怀疑 prompt 约束不足或检索上下文质量差；如果 Answer Relevancy 低，则更多是指令聚焦不够、问题重写有偏差，或者回答模板没有把用户目标钉死。真正成熟的回答不是背指标定义，而是能顺着指标讲出下一步怎么改。
\end{analysisBox}

\begin{figure}[htbp]
  \centering
  \includegraphics[width=0.84\textwidth]{figures/rag-iteration-loop.jpg}
  \caption{离线快速迭代、线上最终验收、两者互相校准}
\end{figure}

\begin{summaryBox}{怎么讲离线和线上闭环}
工程上不能只做离线 benchmark，也不能只盯线上点击率。更合理的做法是：先在离线测试集上跑 Hit@K + RAGAs，定位检索层还是生成层出了问题；修完后再看线上点击率、转人工率、用户满意度有没有一起改善；如果线上暴露了新的坏样本，再把这些样本补回离线测试集，进入下一轮。这样离线是诊断工具，线上是最终验收，两者互相校准。
\end{summaryBox}

\begin{figure}[htbp]
  \centering
  \includegraphics[width=0.9\textwidth]{figures/rag-metrics-table.png}
  \caption{RAG 指标总表：指标所属层级、衡量对象与低分时的优先排查方向}
\end{figure}

\subsection{Embedding 怎么评估，光看 MTEB 榜单行不行}

\begin{summaryBox}{这一题面试官在试什么}
表面问的是“你们 RAG 用什么 embedding、怎么评估好不好”，实际考的是：你能不能把“把问题和文档变成向量”这一步放回真实业务里判断，而不是报一个公开榜单名次就完事。会背 MTEB 和真正做过选型，差距就在这一题。
\end{summaryBox}

\begin{summaryBox}{面试里这题的标准说法}
先拿 50 个真实用户问题，再为每个问题标出“标准应该找到哪份资料”（gold doc）。然后看系统前 10 个候选里有没有这份资料、它排在第几位，专业点说就是用业务 query + gold doc 跑 Hit@K / Recall@K 和 MRR。我的判断是：没有自家测试集的 embedding 选型，本质上只是换模型抽盲盒。MTEB / C-MTEB 只能做第一轮粗筛，把候选从几十个砍到三五个；最终选择必须由业务测试集定。
\end{summaryBox}

\begin{figure}[htbp]
\centering
\tikzset{
  evstep/.style={rectangle, rounded corners, draw=darkblue, line width=0.8pt,
             fill=blue!4, text width=2.55cm, minimum height=1.9cm,
             align=center, inner sep=4pt},
  evnum/.style={circle, fill=orange!85!black, text=white, font=\bfseries\footnotesize,
            inner sep=1.4pt, minimum size=1.4em},
  evarr/.style={-{Stealth[length=2.2mm]}, line width=0.9pt, draw=orange!85!black},
}
\begin{tikzpicture}[font=\small]
\node[evstep] (s1) {业务 query\\$+$\\gold doc 绑定};
\node[evstep, right=1.0cm of s1] (s2) {三件套核心指标\\Hit@K / Recall@K\\/ MRR};
\node[evstep, right=1.0cm of s2] (s3) {五类 query 分组\\同义 / 缩写 / 跨语言\\/ 数字 / 长句};
\node[evstep, right=1.0cm of s3] (s4) {失败样例 trace\\Hit\,=\,0 单独捞出\\定位病灶};
\node[evnum] at (s1.north west) {1};
\node[evnum] at (s2.north west) {2};
\node[evnum] at (s3.north west) {3};
\node[evnum] at (s4.north west) {4};
\draw[evarr] (s1) -- (s2);
\draw[evarr] (s2) -- (s3);
\draw[evarr] (s3) -- (s4);
\end{tikzpicture}
\caption{最小 embedding eval 的四步：绑定 gold doc、跑三件套指标、按五类 query 分组、再对失败样例做 trace}
\end{figure}

\begin{analysisBox}{为什么不能只看榜单，也不能只“换更大的模型”}
MTEB 测的是“通用考试题”，而你的系统面对的是“公司自己的题”：内部工单、合同、代码库、客服记录。公开榜单靠前的模型换到法务 PDF、内部缩写场景上，Top-10 命中率反而可能掉十几个点。“换更大的模型”同样靠不住——不同领域上 text-embedding-3-large、Voyage、bge-m3 这类模型的胜负关系是交叉的，把“大等于准”当标准，恰好跳过了最关键的一步：它在你的资料库里到底找不找得到正确资料。

真正的分水岭不是你报出哪个模型名，而是你能不能拿出一张“哪些问题找得到、哪些找不到”的失败样例表。判断顺序我会固定成三步：先看找不找得到（Hit@K / Recall@K），再看排得靠不靠前（MRR），最后看失败集中在哪类问题上。不要只看 cosine 均值，它更像平均相似度，不能直接说明用户能不能拿到正确资料。
\end{analysisBox}

\begin{qaBox}{单个标准答案和多个标准答案，指标怎么选}
单个 gold doc 的场景先看 Hit@K：前 K 个候选里有没有命中那条正确资料；多个 gold doc 的场景再看 Recall@K：正确资料被找回了多少比例。所以工程上我会把 Recall@K 和 Hit@K 一起记录，而不是只留一个。代码里 \term{record\_retrieval\_eval} 对每条样本同时算 hit、recall@k、reciprocal\_rank，正是为了支撑这个判断顺序。
\end{qaBox}

\subsubsection*{判断核心：按 query 类型分组，比平均分更值钱}

\begin{analysisBox}{为什么必须按 query 类型分组}
平均分会把问题盖住。企业 RAG 最容易翻车的不是标准问法，而是五类“麻烦问题”：同义改写、公司内部缩写与术语、中英混着问的跨语言、订单号/错误码这类精确字符、以及很长很口语的长 query。这些在公开榜单里不一定多，但在你的知识库里天天出现。

工程动作就是：给每条问题打一个类型标签，跑完后按标签分组看 Recall@10；再把失败问题单独拿出来，看是资料切得不对、模型不认识内部词，还是 gold doc 本身标错了。这套数据沉淀下来就是回归集，每次换模型、换切分策略都重跑。代码里 \term{RAGEvalExample.query\_type} 把类型带进每条样本，collector 的 \term{by\_query\_type} 做分组聚合，\term{failures} 把 Hit 为 0 的样本单独 trace 出来，正是这套方法的落地。
\end{analysisBox}

跑完后不要只看平均分，按类型聚合，弱在哪一眼就能看出来。下面是一份典型的分组诊断表，结论不是“换更大的 embedding”，而是分类型对症下药：

\begin{center}
\small
\begin{tabularx}{\textwidth}{>{\bfseries}lllX}
\toprule
Query 类型 & Recall@10 & MRR & 诊断方向 \\
\midrule
同义改写 & 0.92 & 0.71 & 先不动 \\
缩写术语 & 0.34 & 0.18 & 加 BM25 关键词召回 \\
跨语言 & 0.61 & 0.42 & 先看失败样例再定 \\
数字代码 & 0.28 & 0.14 & 关键词精确兜底 \\
长 query & 0.78 & 0.55 & 加 rerank \\
\bottomrule
\end{tabularx}
\end{center}

\begin{figure}[htbp]
\centering
\begin{tikzpicture}[font=\small]
  % axis scale: 1.0 recall -> 5cm
  \def\H{5}
  \def\bw{0.42} % bar width
  \def\gp{2.2}  % group pitch
  % y axis + gridlines
  \draw[->] (-0.4,0) -- (0,0) -- (0,\H+0.5);
  \foreach \y in {0,0.2,0.4,0.6,0.8,1.0}{
    \draw[gray!30] (0,\y*\H) -- (5*\gp+0.2,\y*\H);
    \node[left, gray!60, font=\scriptsize] at (-0.05,\y*\H) {\y};
  }
  \node[rotate=90, anchor=south, font=\scriptsize] at (-0.7,\H/2) {Recall@10};
  % groups: name / baseline / hybrid
  \foreach \i/\name/\base/\hyb in {
    0/同义改写/0.72/0.76,
    1/缩写术语/0.34/0.82,
    2/跨语言/0.68/0.71,
    3/数字代码/0.28/0.75,
    4/长 query/0.61/0.64}{
    \pgfmathsetmacro\x{\i*\gp+0.5}
    % baseline bar (gray)
    \filldraw[fill=gray!45, draw=gray!60]
      (\x,0) rectangle (\x+\bw,\base*\H);
    \node[above, font=\scriptsize, gray!55] at (\x+\bw/2,\base*\H) {\base};
    % hybrid bar (orange)
    \filldraw[fill=orange!85!black, draw=orange!90!black]
      (\x+\bw+0.08,0) rectangle (\x+2*\bw+0.08,\hyb*\H);
    \node[above, font=\scriptsize, orange!85!black] at (\x+1.5*\bw+0.08,\hyb*\H) {\hyb};
    \node[below, font=\scriptsize] at (\x+\bw+0.04,0) {\name};
  }
  % legend
  \filldraw[fill=gray!45, draw=gray!60] (0.3,\H+0.15) rectangle (0.6,\H+0.45);
  \node[right, font=\scriptsize] at (0.65,\H+0.3) {基线（仅换 embedding）};
  \filldraw[fill=orange!85!black, draw=orange!90!black] (4.2,\H+0.15) rectangle (4.5,\H+0.45);
  \node[right, font=\scriptsize] at (4.55,\H+0.3) {混合方案（$+$ rerank）};
\end{tikzpicture}
\caption{仅换 embedding vs hybrid$+$rerank：按 query 类型分组的 Recall@10 对比。缩写术语与数字代码两类收益最大（$+0.48$、$+0.47$），同义改写几乎无差别——印证“分类型对症下药”而非笼统换大模型。}
\end{figure}

\begin{summaryBox}{为什么术语和数字代码要走 hybrid，而不只是换 embedding}
内部缩写、产品型号、订单号不像自然语言，更像精确字符，比如 SKU-X1932、ORD20260418，纯语义检索反而不可靠，关键词检索更稳。更值得做的是 hybrid：语义搜索找一批、关键词搜索找一批，合并后再 rerank。一份匿名复盘可以印证方向：仅换 embedding，整体 Recall@10 从 0.74 到 0.78；加关键词兜底 + 重排后整体 0.74 到 0.91，缩写类从 0.34 拉到 0.82。这不是公开基准，只说明诊断方向——收益最大的恰恰是缩写类和数字代码类。
\end{summaryBox}

\begin{qaBox}{追问：Recall@10 高但答案仍然差，是 embedding 的问题吗}
大概率不是。前 10 个候选里已有正确资料，说明“找资料”这一步不算坏；问题多半在后面：rerank 没把它推到前三、提示词没要求优先看高排名片段、或引用被裁断。处理顺序是先看 rerank，再调提示词，最后才动 embedding。这也对应生成层诊断图里的结论：Context Recall 不低却 Faithfulness/Answer Relevancy 低，病灶在排序与生成约束，而非召回。
\end{qaBox}

\begin{qaBox}{追问：怎么用 50 条 query 做一个最小 embedding eval}
五步就够：抽 50 条真实问题；给每题标出应该找到的 1\,-\,3 篇 gold doc；跑两个候选模型；看前 5 / 前 10 是否命中、排第几；最后按问题类型分组诊断。50 条不是上限而是起跑线，跑通后再扩到 200。关键是先把评估闭环跑通——同一份 200 条样本上跑过 3 个 embedding，比只看 MTEB 排行榜有用得多。
\end{qaBox}

\subsection{生产追问：Agent Eval 怎么拆到可定位故障}
\begin{qaBox}{工程结论}
业务 Agent 评估不能只给最终答案打分，而要拆成“任务完成、证据召回、工具执行、答案可信、成本受控、恢复能力”六层。这样一旦分数下降，能定位是 retrieval、rerank、generation、tool policy、budget breaker 还是 harness recovery 的问题。
\end{qaBox}

\begin{analysisBox}{按流程讲最清楚}
第一步先定义评估单元。不是所有任务都该用同一把尺子，至少要区分事实问答、研究总结、流程查询、开放分析这几类任务，并给出 gold answer、gold doc 或最小验收标准。第二步跑离线基准。对 RAG 问题先测检索层的 Hit@K / Recall@K / MRR，再测生成层的 Faithfulness / Answer Relevancy / Context Recall / Context Precision。第三步看执行层。也就是工具调用次数、成功率、失败类别、是否触发审批、是否被 budget breaker 截断。第四步再接线上指标，例如点击率、采纳率、转人工率、用户追问率。

每个指标最好直接讲语义。Hit@K 是“找没找到”；MRR 是“找得靠不靠前”；Faithfulness 是“有没有按证据说”；Answer Relevancy 是“有没有答到问题”；Context Recall 是“该给模型看的关键上下文有没有进来”；Context Precision 是“塞给模型的上下文里噪音多不多”。再往执行侧，tool success rate 是“工具面是否稳定”，approval rate 是“高风险动作触发频率”，budget overrun rate 是“系统是否经常失控”。

技术选型上也有两种路。方案 A 是只做 LLM-as-a-judge，优点是搭得快，缺点是会把检索、排序、生成、工具问题混成一锅。方案 B 是分层评估，把 retrieval、generation、execution、business acceptance 拆开。DeepIntel 已经明显走向方案 B，因为项目里既有 \term{record\_retrieval\_eval}，也有 tool history、node outcome、budget state 这些运行时观测面。
\end{analysisBox}

\subsection{生产追问：无增益 tool call 怎么量化和治理}
\begin{qaBox}{工程结论}
这类问题不能只看 answer quality，因为答案可能没变，但效率已经变差了。最直接的做法是补一组“过程效率指标”，例如 \term{tool efficiency}、\term{redundant tool call rate} 和 \term{marginal utility per call}。
\end{qaBox}

\begin{analysisBox}{推荐的评价口径}
至少有两种口径。方案 A 是只看最终答案对不对。优点是简单，缺点是会纵容“结果没坏但过程很浪费”的行为。方案 B 是把结果指标和过程指标拆开：结果层继续看正确率、Faithfulness、任务完成率；过程层新增冗余调用率、平均每次成功任务的 tool calls、同类查询重复调用率，以及“去掉某次调用后结果是否变化”的边际贡献评估。这里显然更该选方案 B，因为 Agent 系统的成本和稳定性本来就是产品质量的一部分。

面试时可以给一个很稳的表述：如果多调用工具但最终答案完全不变，我会把它定义成“无增益调用”。评估上不改最终正确率分，但会拉低过程效率分。更细一点，可以把一次任务的总分写成：\term{quality score} 和 \term{efficiency penalty} 分开，前者奖励答对，后者惩罚多余调用、重复检索和无效重试。这样既不误伤答案质量，也不会让资源浪费被平均分盖住。
\end{analysisBox}

\subsection{生产追问：Chunking 如何避免破坏证据边界}
\begin{qaBox}{项目内实现}
DeepIntel 当前采用的是偏工程稳妥的段落优先切分：先按空行切段，尽量把相邻段落拼到 \term{chunk\_size} 上限内；如果单段过长，再按字符窗口切开，并保留 \term{chunk\_overlap} 重叠。
\end{qaBox}

\begin{analysisBox}{为什么这么做}
当前仓库里的实现比较明确：先把文本做换行归一，再按段落分块；短段落尽量合并，长段落才退化到窗口切分。这种方案的优点是实现简单、稳定、对 markdown / txt / json 转文本后的内容都能工作；缺点是它主要按表面结构切，不理解标题层级、表格语义和代码块边界。

至少有两种方案。方案 A 是固定窗口切分，优点是简单、吞吐高，缺点是容易把语义句子和段落拦腰切断。方案 B 是结构感知切分，例如先按标题、列表、表格、章节切，再在局部做窗口化。优点是语义完整度更高，缺点是要先把文档标准化并做更复杂解析。DeepIntel 当前实际落地更接近 A 和 B 的中间态：不是裸滑窗，而是“段落优先 + 超长回退滑窗”。
\end{analysisBox}

\subsection{生产追问：多格式文档如何做 canonical representation}
\begin{qaBox}{结论}
先不要急着切 chunk，先做文档标准化。核心原则是把不同来源先转成统一的中间表示，再按文档类型决定切分策略，而不是所有输入都直接按字符切。
\end{qaBox}

\begin{analysisBox}{处理顺序}
更稳的设计至少分两层。第一层是解析归一化，把 pdf、docx、json、txt 先抽成统一文本或中间结构；项目里现在已经支持 txt、md、json、docx、pdf，并且 json 会递归展开成可读文本。第二层是类型化切分。制度文档适合按标题和条款切，FAQ 适合按问答对切，代码文档适合按函数 / 类切，工单记录适合按时间段或会话轮次切。

方案 A 是所有内容都 flatten 成纯文本，再统一 chunk。优点是实现最快，缺点是结构信息损失大。方案 B 是先做 canonical representation，例如 \term{title / section / body / table / metadata} 这样的中间层，再按类型路由切分器。长期一定更应该走方案 B，因为“先统一表示，再分类型处理”才有机会把不同文档类别都做好。当前项目已经做了第一步的文本归一化，但严格说，按文档类型分策略这一步还没有完全展开，面试里要诚实讲成“现阶段边界”。
\end{analysisBox}

\subsection{生产追问：Embedding、稠密召回和 Rerank 的取舍边界}
\begin{qaBox}{Embedding 选型原则}
Embedding 选型我会固定讲四件事：语言覆盖、领域匹配、向量维度与成本、和现有检索链路的兼容性。公开榜单只做粗筛，最终一定回到业务测试集。
\end{qaBox}

\begin{analysisBox}{选型对比}
至少有两个方向。方案 A 是选通用大模型 embedding，优点是多语言和泛化通常更强，缺点是成本高、维度大、在内部术语场景未必稳赢。方案 B 是选开源可本地化模型，例如项目里现在配置的 BAAI 系列，优点是可控、成本低、便于私有化，缺点是要自己做业务评估和运维。DeepIntel 当前选的是方案 B，仓库里默认 embedding 是 \term{BAAI/bge-zh-qwen2-int8}，rerank 是 \term{BAAI/bge-reranker-v2-m3}，这是偏中文和可控性的取舍。

判断标准不要只背模型名。语言覆盖看中英混查和多语言资料；领域匹配看内部术语、缩写、数字代码；工程成本看维度、吞吐、部署资源；链路兼容性看它和 BM25、rerank、pgvector 的组合效果。这才是成熟回答。
\end{analysisBox}

\begin{qaBox}{稠密召回的失效模式是什么}
稠密检索的核心是把 query 和 document 映射到同一个向量空间里，再按 cosine similarity 或 inner product 做近邻搜索。它的强项是语义改写，弱项是精确编号、罕见术语、权限过滤和结构化字段，所以生产 RAG 通常要 hybrid search、metadata filter 和 rerank 共同兜底。
\end{qaBox}

\begin{analysisBox}{原理版回答}
如果面试官继续追问，可以这样展开：训练阶段希望相关 query-doc 对在向量空间更近，不相关对更远；推理阶段把 query 编成向量后，在向量库里做近邻搜索，拿回 Top-K 候选。它擅长语义改写、同义表达和自然语言问题，但对精确字符串、编号、缩写不一定稳，所以企业场景常常要和 BM25 做 hybrid。
\end{analysisBox}

\begin{qaBox}{Rerank 为什么不能替代召回}
Rerank 是第二阶段排序，只能重排已召回候选，不能补回第一阶段没找到的证据。因此如果 Recall@K 本身很低，换更强 reranker 不会根治；先要修 chunking、hybrid recall、query rewrite 和 metadata filter。
\end{qaBox}

\begin{analysisBox}{怎么讲算法原理}
至少有两大类。方案 A 是 bi-encoder 召回，也就是 embedding 检索，优点是快，缺点是相关性判断较粗。方案 B 是 cross-encoder rerank，把 query 和 candidate 一起送进模型做交互打分，优点是排序更准，缺点是慢。工程上最常见的组合就是“bi-encoder 召回 + cross-encoder 重排”，DeepIntel 现在就是这个方向，所以它前面要先做 Top-K 截断，后面再把更贵的 reranker 用在少量候选上。
\end{analysisBox}

\subsection{生产追问：RAG 事故如何定位到 retrieval、ranking 还是 generation}
\begin{qaBox}{固定判断顺序}
我会先看“正确 chunk 有没有进 Top-K”，再看“有没有被排到足够靠前”，最后才看“模型有没有按证据说”。也就是先查 retrieval，再查 ranking，再查 generation。
\end{qaBox}

\begin{analysisBox}{诊断路径}
如果 Hit@K / Recall@K 很低，大概率是检索层问题，优先怀疑 chunking、embedding、BM25 兜底、多路召回。若 Hit@K 不低但 MRR 很低，说明“找到了但排得不前”，优先查 rerank、RRF 融合和 Top-K 截断。若 Context Recall 不低，但 Faithfulness 或 Answer Relevancy 低，就更像生成层问题，通常是提示词约束、引用组织、回答模板或模型输出跑偏。

这里也有两个常见误区。第一是看见答案错就立刻换模型；第二是看见召回高就以为检索没问题。更成熟的判断是：召回、排序、生成分层排查。DeepIntel 的离线评估代码本身就是按这个思路拆的，检索层和生成层分开记录，而不是只留一个总分。
\end{analysisBox}

\subsection{trace 为什么要分成几条线}
\begin{itemize}[leftmargin=2em]
  \item \term{agent\_trace}：高层流程事件，适合 UI 和时间线回放。
  \item \term{guardrail\_trace}：安全和路由决策，适合解释“为什么被拦”。
  \item \term{tool\_call\_audit}：单次工具调用明细，适合法证级追踪。
\end{itemize}

\begin{summaryBox}{设计观点}
把所有东西都塞进一条 trace 里，看起来简单，实际上后面既不好查，也不好给不同角色消费。高层事件、护栏决策、单次调用明细这三类数据天然就是不同层次，拆开才有意义。
\end{summaryBox}

\begin{analysisBox}{深入分析}
很多系统刚开始都会把“可观测性”理解成“有日志就行”，再进一步可能是“有一条 timeline 就行”。但当系统同时涉及业务流程、治理决策和外部工具时，单一时间线很快就会失效。原因不是数据量大，而是语义层次不同。产品和前端更关心用户可见的高层进度，例如 Planner 开始了、Report 生成了、任务完成了；安全和治理更关心“为什么这个动作被拦了”“为什么这个节点进入 awaiting\_approval”；法证和排障则需要知道单次工具调用的参数、结果摘要、错误类别和耗时。

如果不拆线，后果通常有两个。第一，查询和展示成本会上升，因为每个消费者都得从同一条混杂日志里自己再切语义层。第二，边界会变模糊，尤其是当你想回答“这个错误到底是业务失败、治理阻断还是外部工具异常”时，很难快速切干净。

所以拆 trace 的核心不是“多存几份”，而是让不同层次的问题有不同层次的数据承载。面试时最好直接说：\term{agent\_trace} 服务于高层流程可见性，\term{guardrail\_trace} 服务于治理解释，\term{tool\_call\_audit} 服务于单次调用法证。这样表达会非常清晰。
\end{analysisBox}

\section{第七轮拷打：哪些能吹，哪些不能吹}

\begin{summaryBox}{这一轮面试官在试什么}
这轮最考验真实性。很多项目不是死在技术上，而是死在表述上。你如果把设计稿说成已经上线，或者把半成品说成生产级，懂行的面试官一追就穿。这里必须把“已做完”“已部分落地”“还在设计”分开讲。
\end{summaryBox}

\begin{qaBox}{这一轮最常见的追问}
\begin{itemize}[leftmargin=2em]
  \item 你现在到底能不能说自己做了治理闭环。
  \item 哪些能力已经在代码里落地，哪些只是详细设计。
  \item 历史上最关键的 bug 是什么，修复后语义有什么变化。
  \item 如果我现在把系统部署到生产，会先卡在哪几个现实问题上。
\end{itemize}
\end{qaBox}

\subsection{当前真实完成度应该怎么诚实表达}
\begin{riskBox}{建议表述}
当前系统已经具备研究主链路、DAG 编排、反思校验、SSE trace 和基本预算治理能力，但完整治理闭环仍在演进中。更准确的说法是：它已经是一个能工作的研究型 Agent 系统，但离严格意义上的生产级执行治理系统还有距离。
\end{riskBox}

\begin{center}
\small
\begin{tabularx}{\textwidth}{>{\bfseries}lXX}
\toprule
能力项 & 当前状态 & 面试表述建议 \\
\midrule
主任务生命周期 & 已打通 & 可以说“已实现主链路” \\
失败语义 & 已从历史 bug 修到显式 outcome & 可以说“失败语义已显式化” \\
工具审计 & 已有独立 \term{tool\_call\_audit} 并改为运行中持续落库 & 可以说“审计已具备 crash-safe 强化版，但仍不是全量法证平台” \\
会话级预算 breaker & 已有硬熔断与中途预算状态持久化 & 可以说“已支持会话级硬熔断，但 token/cost 仍部分依赖估算” \\
动作级审批 & 已接入独立 reviewer 分支与 approval API & 可以说“已实现动作级审批主链路，但恢复粒度仍以批次边界为主” \\
Harness supervisor & 已有图外 supervisor 文件协议与恢复快照 & 可以说“已实现最小可用任务级恢复控制，但还不是多 worker 生产调度器” \\
MCP 安全接入 & 已有 registry + policy proxy + allowlist 闸门 & 可以说“已接入安全通路，但默认业务图还未大规模使用 MCP 节点” \\
\bottomrule
\end{tabularx}
\end{center}

\subsection{失败语义为什么是治理第一优先级}
\begin{center}
\input{figures/failure-routing}
\end{center}

\begin{summaryBox}{历史 bug 与当前状态}
这个系统历史上最关键的问题，是工具节点失败后聚合器可能仍然把节点当成完成。这个 bug 现在已经修掉了，工具节点通过 \term{NodeOutcome} 回传显式状态，聚合器按 \term{success / retryable\_error / terminal\_error / awaiting\_approval / skipped} 更新节点。它之所以值得在面试里重点讲，是因为它直接体现了“为什么失败语义必须显式建模”。
\end{summaryBox}

\begin{qaBox}{原因}
为什么失败语义是治理第一优先级，而不是审批、预算或者 MCP？因为只要失败语义不对，后面所有治理都建立在错的状态之上。审批系统的前提是你知道当前动作到底是在成功、失败、等待审批还是已经终止；预算系统的前提是你知道当前节点是否还应该继续消耗资源；恢复系统的前提则是你知道上一次中断时到底停在了哪种失败态。历史上那个“工具失败但节点被标 done”的 bug，本质上就是完成态被污染。
\end{qaBox}

\begin{qaBox}{选型}
这里至少有两个方案。方案 A 是继续把失败信息藏在 tool history 或异常文本里，由聚合器自己猜当前节点到底算成功还是失败。优点是改动少，缺点是语义脆弱。方案 B 是让工具层显式回传一份最小控制契约，也就是 \term{NodeOutcome}，把成功、可重试失败、终止失败、等待审批和跳过这些状态直接告诉调度层。这里必须选方案 B，因为控制不能建立在猜测审计日志的基础上。
\end{qaBox}

\begin{qaBox}{如何做}
当前实现里，每个工具节点执行后都会显式回传 \term{NodeOutcome}，其中至少包含 \term{success / retryable\_error / terminal\_error / awaiting\_approval / skipped} 这些控制语义，以及错误类别、重试次数、工具调用 ID 等最小信息。聚合器不再根据异常文本猜测结果，而是按这份显式契约更新 DAG 节点状态和 runtime status。与此同时，Browser 的本地 fallback 也不再被上层误记成成功页面，而是会把真实错误向上转换成 retryable 或 terminal 失败。
\end{qaBox}

\begin{qaBox}{效果}
这样做以后，失败语义从审计层附带信息升级成了调度层硬约束。工具失败不再会被吞掉，后续路由、预算治理、审批暂停和恢复决策都有了可信状态基础。面试时可以很明确地说：失败语义已经显式化，这个历史 bug 已经修掉，不应继续被当成现存问题对外表述。
\end{qaBox}

\subsection{会话级预算治理现在到了什么程度}
\begin{qaBox}{原因}
预算治理最容易被答成“有个 max\_tokens 就行”，这其实远远不够。研究型 Agent 的成本不是单一来源，至少包括模型 token、外部工具调用、网页抓取时长、重试次数和整体 wall clock。只盯 token，会漏掉 Browser 一直超时、MCP 工具反复重试、或者 Planner 把图拆得过深这类真实成本问题。
\end{qaBox}

\begin{qaBox}{选型}
这里至少有两个方案。方案 A 是 observability-only，只打 warning，不改系统行为。优点是实现最简单，对现有流程侵入小；缺点是出了预算事故你只能复盘，不能止损。方案 B 是 breaker 模式，让预算阈值真正参与路由决策。比如超过某个阈值后停掉可选 Web 扩展、禁止新的 replan、必要时提前进入部分报告。这里必须选方案 B，因为治理的定义就是让规则能改变执行，而不只是描述执行。
\end{qaBox}

\begin{qaBox}{如何做}
当前项目已经不是只有长度档位了，图内已经接入 \term{session-level budget breaker}，包括 \term{max\_tool\_calls}、\term{max\_wall\_clock\_seconds}，以及 \term{max\_total\_tokens / max\_cost\_usd} 的运行时累计通道。预算治理分三层。第一层是计量层，尽可能记录 usage、cost、tool call 和耗时；拿不到精确值时明确标估算，并把 \term{estimated} 与 \term{usage\_source} 写入审计。第二层是策略层，在 80\% 左右给 warning，在硬阈值时停止可选扩张，在更高风险时禁止 replan。第三层是收口层，即使被硬熔断，也不会让任务直接烂尾，而是跳过剩余工具节点，进入分析与部分报告，同时把预算中间态持续落到 \term{session\_budget\_state}。
\end{qaBox}

\begin{qaBox}{效果}
这样做以后，预算治理不再只是一个提醒系统，而是真正能改变执行路径的硬约束。系统会在预算接近危险区时发出告警，在硬阈值时主动止损并优雅降级。这里也必须诚实说明边界：当前 token/cost 通道虽然已经接到 LLM usage 和 budget state，但它仍不是绝对精确计费系统，因为上游 provider 不总返回 usage，tool 侧也未必暴露真实账单，所以系统只能做到可精确则精确，不可精确则显式标记估算。
\end{qaBox}

\subsection{法证级审计现在到了什么程度}
\begin{qaBox}{当前结论}
现在已经不只是 \term{tool\_histories} 运行时累积状态了，系统已新增独立的 \term{tool\_call\_audit} 表，把单次工具调用拆成独立记录，并在 workflow 运行过程中持续 upsert。这样 session 结果查询和审计查询不再混在一行 JSON 里，中途崩溃时也不会像以前那样整批丢光。
\end{qaBox}

\begin{riskBox}{仍未闭环的点}
现在工具审计和预算状态已经前移到运行中持续落库，所以 crash-safe 能力比之前强很多。但它仍不是“无限细粒度恢复”：当前最可靠的恢复语义仍是批次边界恢复，不承诺单节点执行到一半的原地续跑。
\end{riskBox}

\begin{analysisBox}{深入分析}
“法证级审计”这四个字不能乱用。真正的法证关注的是：系统在某个时刻为什么做了某个动作、当时依据是什么、参数是什么、结果是什么、谁批准的、是否经过脱敏、崩溃前最后一条可确认事实是什么。单靠最终 session JSON 或一条高层 trace，是达不到这个要求的。

当前系统之所以比早期状态前进了一大步，是因为工具调用和预算状态已经从“流程结束后统一汇总”前移成“运行中持续落库”。这意味着进程中途崩掉时，至少已经完成的调用明细、预算累计和关键状态切换不会整批蒸发。这就是 crash-safe 的核心价值，不是恢复得有多优雅，而是证据先别丢。

但也必须守住边界。现在的强项是法证记录比以前完整得多，不是说恢复已经精确到单条工具流内部。批次边界恢复依然是主语义，所以更准确的表述应该是“已经具备 crash-safe 的实时审计闭环雏形，并支持批次边界恢复”，而不是“已经完成全量执行回放平台”。这类表述差异，面试里非常关键。
\end{analysisBox}

\section{第八轮拷打：邪恶 Agent 怎么管住}

\begin{summaryBox}{这一轮面试官在试什么}
这轮不是问“你会不会做 Agent”，而是问“你敢不敢把 Agent 放进真实系统”。如果面试官开始追问审批、预算、审计、MCP、恢复、权限边界，说明他在看你有没有治理视角，而不只是能力视角。
\end{summaryBox}

\begin{qaBox}{这一轮最常见的追问}
\begin{itemize}[leftmargin=2em]
  \item 失败后到底是重试还是终止，谁来裁决。
  \item 高风险动作要不要审批，审批状态到底谁说了算。
  \item tool audit 能不能支撑事后法证，而不只是页面展示。
  \item 如果明天接 MCP 工具，为什么不能直接把它们当普通 Tool 放进图里。
\end{itemize}
\end{qaBox}

\subsection{治理闭环现状判断}
\begin{summaryBox}{一句话}
现在系统已经有 workflow、trace、tool audit、budget breaker、动作级审批、Harness supervisor 和 MCP policy proxy 的主干治理链路，但恢复粒度、MCP 默认用量和计费精度仍有边界。
\end{summaryBox}

\subsection{审批状态源（Approval Source of Truth）}
\begin{center}
\input{figures/approval-sot}
\end{center}

\begin{qaBox}{固定立场}
动作级审批的最终权威源必须是本地数据库审批单。Harness、Slack、Jira、ServiceNow 这类系统都可以作为审批信号源，但最终 \term{can\_execute(action)} 的判断必须由本地审批聚合根负责。
\end{qaBox}

\begin{analysisBox}{深入分析}
这道题背后其实是在问“source of truth 到底怎么定”。很多团队一开始会误以为，用现成平台做审批最省事，所以直接把外部状态当执行依据就行。但一旦动作进入自动化执行流，审批不再只是一个页面状态，而是一个会影响幂等、恢复、重试和危险动作下发的强控制条件。这时如果没有本地单一真相源，系统几乎必然会出现状态漂移。

为什么数据库审批聚合根最适合做最终裁决？因为动作执行发生在你的系统里，执行前最后一跳必须依赖一个你能稳定查询、能事务更新、能做版本控制、能和审计关联的本地状态。如果直接去问 Harness、Slack 或 Jira 当前是否 approved，你会立刻遇到几个问题：远端状态变化有延迟，回调可能丢失，远端 UI 和本地执行记录不在同一事务里，失败重试时也缺统一幂等点。

所以最稳的架构不是拒绝外部审批，而是重新定义它们的角色：外部系统负责承载人类交互和审批信号，本地数据库负责承载最终执行裁决。只要把这层权责讲清楚，面试官通常会认可你不是在排斥外部平台，而是在保护动作执行的一致性边界。
\end{analysisBox}

\begin{center}
\small
\begin{tabularx}{\textwidth}{>{\bfseries}lXXX}
\toprule
来源 & 谁是权威 & 优点 & 问题 \\
\midrule
LOCAL DB & 系统自己 & 单一真相源、查询直观、审计强 & 工程量最大 \\
HARNESS & Harness 平台 & UI、RBAC、流程现成 & 易和业务状态漂移 \\
EXTERNAL & Slack/Jira/OA 等 & 贴合企业习惯 & 远端状态和本地状态同步复杂 \\
\bottomrule
\end{tabularx}
\end{center}

\begin{lstlisting}[caption={ApprovalAggregate 的建议最小字段}]
approval_id
action_id
approval_source
external_id
status
version
approved_by
expire_time
\end{lstlisting}

\subsection{为什么外部审批只能是信号，不能是唯一裁决}
\begin{qaBox}{原因}
因为只要允许多个 source of truth 并存，就一定会遇到状态漂移：例如 Harness 已经批了，但 DB 仍是 pending；或者 Slack 上已经通过，但 agent 因为没看到本地状态变化再次发起执行。真正要避免的不是“查不到谁批了”，而是“重复执行危险动作”。
\end{qaBox}

\subsection{Harness 外挂设计}
\begin{center}
\input{figures/harness-supervisor}
\end{center}

\begin{qaBox}{架构答案}
Harness 应该放在 LangGraph 外层。LangGraph 负责编排业务步骤，Harness 负责任务租约、失败恢复、暂停恢复、预算调度和 supervisor 语义。数据库负责审批、审计和预算事实，Harness 负责 lease、checkpoint 选择和 worker 协调。当前仓库已经落了最小可用实现：有 \term{harness-tasks.json}、progress log 和 state snapshot，但还不是完整的分布式 worker 调度系统。
\end{qaBox}

\begin{analysisBox}{深入分析}
Harness 设计的关键不是“多一份文件”，而是明确任务级控制面和业务执行面的解耦边界。一个长任务系统真正难的，不是把 DAG 跑完，而是在失败、暂停、恢复、多 worker、审批等待这些现实条件下仍然能保持状态一致。Harness 文件之所以重要，是因为它把这些任务级信息从短生命周期的 HTTP 请求和进程内存里抽了出来，变成跨会话可恢复的外部上下文。

这也是为什么 progress file 和 tasks file 在这套设计里不是附件，而是恢复上下文本身。任务一旦跨会话执行，agent 的上下文窗口不再可靠，单靠口头说明或者控制台日志也不可靠。真正能支撑恢复的，是结构化任务状态、最近 checkpoint、当前 batch、审批等待点、预算使用量这些可以被下一次执行直接消费的数据。

当然，这里也要诚实讲代价。外置 Harness 以后，一定会引入双写和状态同步问题，所以必须分域定权：审批、预算、审计以 DB 为准；worker lease、checkpoint 选择、恢复编排以 Harness 为准。只要把这个边界说清楚，面试官反而会认为你想清楚了复杂度，而不是机械套了一个 supervisor 名词。
\end{analysisBox}

\subsection{MCP 安全接入原则}
\begin{center}
\input{figures/mcp-proxy}
\end{center}

\begin{qaBox}{固定表述}
MCP 是协议，不是治理系统。真正安全的接法不是把 MCP Tool 直接映射成 LangChain Tool，而是先过 \term{McpPolicyProxy}，把 server allowlist、tool allowlist、只读默认和 secret redaction 做成强闸门。当前仓库里这条代理通路和 \term{mcp\_server\_registry} 已经落地，但 Planner 还不会默认大量生成 MCP 节点，所以它更适合表述成“安全接入通路已存在”，而不是“业务主流量已全部走 MCP”。
\end{qaBox}

\begin{analysisBox}{深入分析}
MCP 的危险点在于它天然会给人一种错觉：协议既然统一了，接上就能安全使用。实际上协议只解决了“如何发现和调用工具”，没有解决“这个工具该不该被信任、这组参数会不会越权、返回结果有没有敏感字段、写操作需不需要审批”。也就是说，MCP 统一的是接口，不是治理。

如果你把 MCP server 暴露出来的工具直接映射成普通 Tool，有三个高风险后果。第一，server 信任边界缺失。你可能根本不知道这个 server 是谁提供的、是否可信、tool 列表有没有变更。第二，能力边界缺失。同一个协议既可以暴露只读查询，也可以暴露写文件、发消息、调内部系统，如果没有 allowlist 和只读默认，agent 会天然尝试更强能力。第三，数据边界缺失。很多 MCP 返回值会带敏感字段、token、内部 ID，如果不在代理层做 redaction，这些数据会直接进入审计和上下文。

所以 policy proxy 的核心价值，是在业务图和外部工具生态之间插入一层强治理边界。它不是为了让调用更麻烦，而是为了把“能调用”变成“在规则内可调用”。这一层如果讲得足够清楚，面试官通常会把你归到真正考虑过安全接入的人，而不是只会接协议的人。
\end{analysisBox}

\subsection{生产追问：MCP 协议层、传输层和治理层怎么拆}
\begin{qaBox}{工程结论}
MCP 底层消息是 \term{JSON-RPC 2.0}，本地常用 \term{stdio}，远程推荐 \term{Streamable HTTP}。工具 schema 要同时关注 \term{inputSchema} 和 \term{outputSchema}：输入决定能不能执行，输出的 \term{structuredContent} 决定能不能进入后续上下文和审计。
\end{qaBox}

\begin{analysisBox}{怎么回答既不空泛也不吹过头}
如果面试官问协议，最容易答错的点就是把“传输方式”和“消息格式”混在一起。更准确的说法是：MCP 不管跑在本地还是远程，底层消息格式统一都是 \term{JSON-RPC 2.0}。它解决的是“client 调 server 的 method，server 返回 result 或 error”这件事，本质上就是一个轻量 RPC 规范。之所以选它，不是因为花哨，而是因为 JSON 可读、跨语言实现简单、调试成本低，Python 和 TypeScript 写的 server 都能讲同一种消息语言。
\end{analysisBox}

\begin{figure}[htbp]
  \centering
  \includegraphics[width=0.9\textwidth]{figures/mcp-json-rpc-format.png}
  \caption{MCP 底层消息格式统一使用 JSON-RPC 2.0：传输方式可变，消息结构不变}
\end{figure}

\begin{analysisBox}{本地传输：stdio}
本地场景下，MCP 最常见的不是网络，而是 \term{stdio}。Client 启动 MCP Server 这个子进程，然后通过标准输入发 JSON-RPC 请求、从标准输出读 JSON-RPC 响应。两个进程在同一台机器上运行，走的是操作系统管道，不经过网卡，也不需要 TCP/IP 协议栈，所以延迟低、配置简单、没有额外的网络暴露面。
\end{analysisBox}

\begin{figure}[htbp]
  \centering
  \includegraphics[width=0.92\textwidth]{figures/mcp-stdio-pipe.png}
  \caption{stdio 本地通信：Client 把 Server 当作子进程启动，通过 stdin/stdout 交换 JSON-RPC 消息}
\end{figure}

\begin{analysisBox}{远程传输：Streamable HTTP}
远程场景下，现在推荐的是 \term{Streamable HTTP}。核心不是 WebSocket，而是单个 HTTP 端点同时承接请求和响应。Client 通过 POST 发送 JSON-RPC 消息；Server 可以直接返回普通 JSON，也可以在需要流式输出时返回 \term{SSE} 流。关键点在于：这里并没有抛弃 SSE，而是把它收编进单端点架构里，变成“需要流式时的返回形式”，而不是早期那种额外维护一条独立 SSE 长连接。
\end{analysisBox}

\begin{figure}[htbp]
  \centering
  \includegraphics[width=0.92\textwidth]{figures/mcp-streamable-http.png}
  \caption{Streamable HTTP：Client 用 POST 发 JSON-RPC，请求返回普通 JSON 或 SSE 流}
\end{figure}

\begin{analysisBox}{为什么早期 SSE 双端点方案被弃用}
这里要再补一个时间线。早期 MCP 远程方案是 \term{HTTP + SSE 双端点}：一条 POST 端点负责发请求，另一条 SSE 长连接负责收流。这套方案在早期规范里存在过，但在后续规范演进里已经被标记为 deprecated，新项目更应该直接用 Streamable HTTP。原因并不玄学，就是双通道让状态管理、排障和部署都更复杂；而 Streamable HTTP 把同一轮交互收拢回一个端点，语义更简单，也更适配负载均衡和 serverless 场景。
\end{analysisBox}

\begin{figure}[htbp]
  \centering
  \includegraphics[width=0.9\textwidth]{figures/mcp-legacy-remote-sse.png}
  \caption{早期 MCP 远端传输：POST 与 SSE 双端点分离，链路更复杂}
\end{figure}

\begin{figure}[htbp]
  \centering
  \includegraphics[width=0.92\textwidth]{figures/mcp-transport-evolution.png}
  \caption{MCP 远端传输演进：从双端点 SSE 过渡到单端点 Streamable HTTP}
\end{figure}

\begin{analysisBox}{schema 这一层应该怎么补}
如果面试官问的是 schema，不要答成“就是一个 JSON”。更准确的说法是：MCP 工具暴露结构化输入规范，本质上接近 JSON Schema，用来声明参数类型、必填项、枚举和对象结构；2025-11-25 规范还把 \term{outputSchema} 和 \term{structuredContent} 的关系讲得更清楚。DeepIntel 本轮已经把 policy proxy 从只校验输入扩展到校验输出结构，坏的 \term{structuredContent} 不再直接进入 summary、hash 和后续 Agent 上下文。
\end{analysisBox}

\begin{analysisBox}{最后再落回 DeepIntel}
这里也至少有两个接法。方案 A 是把 MCP 当成“统一工具接口”，接上就直接让 Agent 调；优点是快，缺点是把协议误当治理。方案 B 是像本项目这样，把 MCP 当成外部能力暴露协议，但前面一定要加 policy proxy，先做 server trust、tool allowlist、只读默认、审批和 redaction。DeepIntel 明显选的是方案 B，所以面试里最稳的表述不是“我们已经全面 MCP 化”，而是“我们已经把安全接入通路搭好了”。
\end{analysisBox}

\subsection{为什么 Agent 调工具时更容易出事}
\begin{figure}[htbp]
  \centering
  \includegraphics[width=0.96\textwidth]{figures/tool-safety-contrast.png}
  \caption{概率性推理 vs 确定性规则：安全不能押在“模型大概率不会出错”上}
\end{figure}

\begin{summaryBox}{核心判断}
工具调用的风险不在于 Agent 会不会说错话，而在于它一旦带着错误参数碰到真实资源，后果会立刻变成工程事故。所以安全不能建立在“模型大概率不会犯错”上，而必须建立在 Harness 工具层的确定性规则上。
\end{summaryBox}

\begin{qaBox}{原因}
面试官问“Agent 调工具，风险到底在哪里”，本质不是在考你会不会调 API，而是在考你有没有把工具调用链路当成高风险执行面来治理。数据库查询、内部检索、发消息、写文件，这些动作一旦从“语言输出”变成“真实执行”，风险类型就从幻觉升级成越权、注入、数据外泄和副作用失控。
\end{qaBox}

\begin{qaBox}{选型}
这里至少有两个方案。方案 A 是把安全主要押在模型理解力上，相信模型足够聪明就不会发危险请求。优点是接入快、系统侵入小；缺点是它本质上是概率控制，遇到 prompt injection、越权诱导、异常路径时没有硬约束。方案 B 是在 Agent 和真实资源之间加一层 Harness 工具层，把每次调用都当成独立请求，串行经过权限控制、参数校验和异常兜底。这里必须选方案 B，因为安全要求是确定性的，不能建立在“模型大部分时候应该没问题”这种判断上。
\end{qaBox}

\begin{analysisBox}{深入分析}
很多系统早期会误以为，只要把 system prompt 写得更严格一点，或者让模型“理解不要乱查数据库”，安全问题就差不多解决了。这种做法能降低一部分误用，但解决不了根本矛盾：模型的推理输出是概率性的，而权限、参数合法性和敏感信息边界是确定性的。两者不是一个层级的问题。

也正因为如此，真正成熟的工具接入层不会假设 Agent 可信，而是把 Agent 当成一个可能犯错、可能被诱导、也可能被脏输入污染的请求发起者。工具层的职责不是“帮 Agent 把请求转发得更漂亮”，而是先判断这个请求到底有没有资格执行、参数是否合法、失败后会不会把内部结构漏出去。把这层说清楚，面试官通常就会意识到你理解的不是 prompt engineering，而是执行面治理。
\end{analysisBox}

\subsection{Harness 工具层怎么挡住越权和带毒参数}
\begin{figure}[htbp]
  \centering
  \includegraphics[width=0.96\textwidth]{figures/tool-guardrail-flow.png}
  \caption{先控权限，再验参数，受控执行，统一兜底}
\end{figure}

\begin{qaBox}{架构答案}
最稳的设计是把 Harness 放在 Agent 和真实工具之间，所有工具请求先进入受控目录，再经过调用前拦截，最后进入沙箱和受限数据库接口。整个链路至少有三道串行防线：权限控制回答“你有没有资格调这个工具”，参数校验回答“你传进来的参数是不是安全合规”，异常兜底回答“执行失败或返回敏感内容时怎么安全收口”。
\end{qaBox}

\begin{qaBox}{选型}
这里也至少有两个方案。方案 A 是只在工具函数内部各自写保护逻辑，让每个工具自己判断权限、自己处理异常。优点是实现快，局部改动小；缺点是规则分散、容易漏判，而且不同工具的行为不一致。方案 B 是把权限、schema、审批、脱敏、审计收拢到统一 Harness 工具层，工具本体只负责业务能力，治理规则由中间层集中裁决。这里更应该选方案 B，因为统一治理才能做到一致审计、集中升级和纵深防御。
\end{qaBox}

\begin{analysisBox}{流程拆解}
整条链路可以概括成八个字：先控权限，再验参数，受控执行，统一兜底。第一步，Agent 发起工具请求，但它表达的只是意图，不应该直接拥有数据库或外部系统访问权。第二步，请求进入 Harness 工具目录，每个工具都要提前声明能力描述、参数 schema、权限级别和是否只读。Agent 能看见哪些工具，由暴露策略决定，而不是想调什么就调什么。

第三步是调用前拦截，也是含金量最高的一步。权限侧要结合身份、角色、租户、会话上下文做综合鉴权；参数侧要做类型、长度、枚举、必填项校验，并且强制数据库访问走参数化模板，禁止自由拼接 SQL。第四步是沙箱执行层。就算前面都通过了，也不能假设工具实现一定安全，所以仍要限制网络、文件系统写入、CPU/内存和执行时长，把副作用约束在小范围里。

第五步是受控数据库接口。这里不应该给工具完整读写权限，而是给只读账号、只读视图和预定义 SQL 模板，并在底层自动追加租户或用户范围过滤条件。第六步是调用后兜底。结果要先统一结构，再做敏感字段脱敏、返回行数限制和错误净化，最后把谁调用了什么工具、查了什么范围、为什么被拦或被放行写入审计日志。这样下游拿到的始终是规范化、安全化的数据，而不是裸结果。
\end{analysisBox}

\begin{qaBox}{两类典型风险怎么回答}
第一类是参数带毒，也就是注入风险。工程判定口径是：先过 schema，不合法参数直接拒绝；数据库访问强制参数化模板，堵死自由拼接；执行异常统一包装，不能把原始 SQL 报错和内部结构直接暴露给模型。第二类是调用越权和数据越界。工程判定口径是：第一层在工具目录和鉴权器卡调用资格，第二层在数据库接口自动追加租户/用户过滤条件。前一层漏判了，后一层仍然能兜住。
\end{qaBox}

\begin{lstlisting}[caption={Harness 工具层统一输出结构示意}]
{
  "status": "success | denied | safe_error",
  "source": "harness_tool_layer",
  "data": { "...": "masked business payload" },
  "meta": {
    "tool_name": "query_transactions",
    "row_count": 20
  },
  "safety": {
    "masked": true,
    "sanitized_error": true,
    "audited": true
  }
}
\end{lstlisting}

\subsection{五类幻觉怎么治理}
\begin{itemize}[leftmargin=2em]
  \item 知识幻觉：先检索，证据不足就拒答，不能拿模型记忆补业务事实。
  \item 工具幻觉：Observation 必须来自真实工具返回，不能让模型伪造 tool result。
  \item 参数幻觉：schema 校验不过就不执行，工具调用是接口，不是聊天。
  \item 证据幻觉：分析和报告必须能回溯到证据和 citation。
  \item 行动幻觉：高风险工具默认不执行，需要动作级审批门。
\end{itemize}

\section{第九轮拷打：高压追问怎么接}

\begin{summaryBox}{这一轮面试官在试什么}
到这一轮，面试官往往已经默认你做过一些东西了。他接下来要看的是：当问题被压得很细、很尖锐时，你能不能既守住事实边界，又把架构判断讲透。下面这些问题适合直接背熟成口头版。
\end{summaryBox}

\begin{riskBox}{高压回答原则}
高压追问时，不要急着追求“说满”。最稳的做法是先给结论，再补理由，最后补边界。也就是先说“我为什么这样做”，再说“备选方案为什么没选”，最后说“哪些部分我现在还没做完”。这样能显著降低被继续追着打穿的概率。
\end{riskBox}

\subsection{为什么选择 LangGraph，而不是直接做多 Agent 对话}
\begin{qaBox}{答案}
多 Agent 对话更适合开放式协作，不适合强流程研究任务。研究任务要求节点、依赖、预算、审计和恢复都可显式表达，LangGraph 这种状态机编排比纯消息流更适合工程化落地。
\end{qaBox}

\subsection{为什么动作审批的权威源必须是 DB}
\begin{qaBox}{答案}
因为动作执行是高风险事实，不能依赖多个远端状态拼凑出来。数据库里的审批聚合根必须是唯一裁决者，否则系统会出现“远端已批、本地未落”“远端已拒、本地重试继续跑”这类危险分叉。
\end{qaBox}

\begin{analysisBox}{深入分析}
这个问题本质上是在问 source of truth。动作审批如果没有单一权威源，系统迟早会出现状态漂移。比如 Harness 那边显示 approved，Slack 消息里也有人点了批准，但本地任务状态还停在 pending；或者外部系统已经拒绝，本地因为网络抖动没收到回调，任务恢复时又把动作继续发了出去。对于读操作，这类漂移已经麻烦；对于写文件、发消息、调用内部系统的高风险动作，这就是事故入口。

至少有两个方案。方案 A 是直接以外部审批系统或 Harness 状态作为最终依据，本地只做被动查询。优点是接入快，缺点是幂等、审计、一致性和恢复都会依赖远端可用性。方案 B 是数据库审批单作为聚合根，外部系统只提供信号，本地统一落库后再由本地状态决定是否 can\_execute。这里必须选方案 B，因为动作执行最终发生在你的系统里，责任也落在你的系统里，裁决权就不能外包。

再往深一点看，DB 作权威源的意义不只是“查起来方便”，而是让审批和执行共享一个可事务化的判断平面。这样才能把 approval\_id、action\_id、approved\_by、expired、decision\_version 这些信息和执行行为绑定起来，形成真正可追责的闭环。面试时把这一层讲明白，会显得你不仅懂工作流，还懂高风险动作系统的责任边界。
\end{analysisBox}

\subsection{为什么外部审批不能直接作为执行依据}
\begin{qaBox}{答案}
外部审批更适合作为输入信号，不适合作为最终执行前判定。标准做法是外部系统回调统一落库，更新本地审批状态，再由本地状态决定是否允许动作执行。这样幂等、审计和重试逻辑才有单一依赖点。
\end{qaBox}

\begin{analysisBox}{深入分析}
这题和上一题相邻，但重点不同。上一题强调谁是最终裁决者，这一题强调为什么远端状态不能直接驱动执行。核心原因是执行系统需要的是稳定、可重放、可幂等的本地判断，而外部审批系统通常只保证“给你一个事件”，不保证这个事件会在你最需要的时候仍然可查询、可事务绑定、可与本地恢复流程完全同步。

如果直接以外部状态做执行前判断，会出现三类风险。第一是时序风险，回调到了但本地还没持久化，worker 就先继续跑了。第二是幂等风险，同一个 approve 事件重放两次，系统没有本地版本控制时可能执行两次。第三是一致性风险，外部系统后续撤销、过期或权限变更，本地没有统一聚合层就很难判断该如何处理已有任务。

所以外部审批最适合的角色是“信号源”，不是“执行开关”。你的系统应该消费这个信号，更新本地审批聚合，然后只相信本地聚合状态。这样一来，重试逻辑只看本地，恢复逻辑只看本地，审计逻辑也只需解释本地状态如何变化。对工程系统来说，这种单点依赖远比每次执行前都去问一圈外部世界稳得多。
\end{analysisBox}

\subsection{为什么要在 LangGraph 外再加 Harness}
\begin{qaBox}{答案}
因为如果把失败恢复、预算熔断、审批暂停、worker 抢租约和任务恢复都塞进图里，图就会从“业务工作流”膨胀成“超级控制器”。Harness 的价值在于把任务级控制面从业务级编排面剥离出去。
\end{qaBox}

\subsection{为什么 MCP 必须先过 policy proxy}
\begin{qaBox}{答案}
因为 MCP 只定义了如何暴露工具，不定义工具是否可信、参数是否越权、结果是否含敏感信息、危险写操作是否需要审批。没有 policy proxy，MCP 接入越多，系统越接近裸奔。
\end{qaBox}

\subsection{预算治理为什么不能只做软提示}
\begin{qaBox}{答案}
软提示只能告诉你“快超了”，不能真正止损。对长任务 Agent 来说，真正需要的是硬熔断加优雅降级：达到阈值后停止可选工具节点，禁止进一步 replan，最后基于当前证据输出部分报告。这才叫治理，不叫提醒。
\end{qaBox}

\begin{analysisBox}{深入分析}
这道题的关键，不在于预算提醒有没有价值，而在于提醒本身不具备控制力。软提示适合人类在控制台里观察，例如“这个任务已经用了 70\% 预算”；但 Agent 是自动执行系统，如果它收到提示却没有强制约束，执行路径不会自动收缩，成本照样继续增长。

从治理角度看，预算问题至少有两种处理模式。方案 A 是 observability-only，只打 warning，不改系统行为。优点是实现最简单，对现有流程侵入小；缺点是出了预算事故你只能复盘，不能止损。方案 B 是 breaker 模式，让预算阈值真正参与路由决策。比如超过某个阈值后停掉可选 Web 扩展、禁止新的 replan、必要时提前进入部分报告。这里显然要选方案 B，因为治理的定义就是让规则能改变执行，而不只是描述执行。

另外还要强调“优雅降级”这半句。硬熔断不是粗暴失败，尤其对研究任务来说，很多时候当前证据已经足够生成一个部分可用的结果。真正成熟的系统不是预算一到就直接报错，而是知道什么时候该停、停了之后怎样最小化价值损失。把这一点讲出来，会让你的预算治理显得完整，而不是只有一个 kill switch。
\end{analysisBox}

\section{附录：Agent 岗底层推理拷打}

\subsection{底层题：不用框架写 multi-head attention，如何证明 shape 和 mask 没错}
\begin{qaBox}{回答策略}
这题不是为了背 Transformer，而是考你能不能把模型底层算子、张量形状、mask 语义和数值稳定性讲清楚。回答时先给单头 scaled dot-product attention，再扩到 multi-head，最后说明 causal mask、padding mask、softmax 稳定化和复杂度。
\end{qaBox}

\begin{analysisBox}{手撕时最稳的讲法}
multi-head attention 的核心不是“复制很多个 attention”，而是把输入先线性投影成多组 Q、K、V，每一组 head 在自己的子空间里算一次 scaled dot-product attention，再把多个 head 的结果拼接起来，最后过一层输出投影。公式可以压成四步：一，$Q=XW_Q,\ K=XW_K,\ V=XW_V$；二，按 head 数把最后一维 reshape 成 \term{[B, H, T, D\_h]}；三，算 $\mathrm{softmax}(QK^T / \sqrt{D_h})V$；四，把所有 head concat 回去再乘 $W_O$。

至少有两个实现口径。方案 A 是纯 Python / 原生列表手写，优点是方便解释流程，缺点是代码长、性能差、容易写错维度。方案 B 是 NumPy 版本，优点是更接近工程实现，也更容易说明矩阵乘法、transpose 和 broadcasting。面试里更建议先给 NumPy 版骨架，再口头说明如果不用 NumPy 就是把矩阵乘法拆成三层循环。
\end{analysisBox}

\begin{lstlisting}[language=Python,caption={NumPy 版 multi-head attention 最小骨架}]
import numpy as np

def softmax(x, axis=-1):
    x = x - np.max(x, axis=axis, keepdims=True)
    e = np.exp(x)
    return e / np.sum(e, axis=axis, keepdims=True)

def multi_head_attention(x, Wq, Wk, Wv, Wo, num_heads, mask=None):
    # x: [B, T, D]
    B, T, D = x.shape
    Dh = D // num_heads

    Q = x @ Wq
    K = x @ Wk
    V = x @ Wv

    Q = Q.reshape(B, T, num_heads, Dh).transpose(0, 2, 1, 3)
    K = K.reshape(B, T, num_heads, Dh).transpose(0, 2, 1, 3)
    V = V.reshape(B, T, num_heads, Dh).transpose(0, 2, 1, 3)

    scores = np.matmul(Q, K.transpose(0, 1, 3, 2)) / np.sqrt(Dh)
    if mask is not None:
        scores = np.where(mask, scores, -1e9)

    attn = softmax(scores, axis=-1)
    out = np.matmul(attn, V)
    out = out.transpose(0, 2, 1, 3).reshape(B, T, D)
    return out @ Wo
\end{lstlisting}

\begin{qaBox}{高级追问}
最常见追问有五个。第一，为什么除以 $\sqrt{D_h}$：维度大时点积方差会变大，不缩放 softmax 容易过饱和。第二，mask 为什么放在 softmax 前：因为要从概率归一化前移除非法位置。第三，padding mask 和 causal mask 如何广播到 \term{[B,H,T,T]}。第四，长上下文复杂度为什么是 $O(T^2)$，工程上怎么用 KV cache、chunk attention 或检索外存缓解。第五，为什么这个底层题和 Agent 岗有关：因为 Agent 长上下文、工具轨迹压缩和证据窗口管理都依赖你理解 attention 的成本与边界。
\end{qaBox}

\input{advanced-agent-notes}
\input{agent-position-cardbook}
\section{生产级韧性与安全重构：设计说明与面试答法}

\begin{summaryBox}{这一章回答什么}
这一章不是把安全名词罗列一遍，而是把项目从“功能能够跑”提升到“外部依赖失败时仍可解释、数据越权时默认拒绝、状态恢复时有事实依据”的设计完整讲清楚。每个主题都按同一套问题展开：为什么需要它，采用什么方案，为什么没有选其他方案，收益是什么，代价是什么，已经考虑了哪些边界，还有哪些边界不能声称已经解决。
\end{summaryBox}

\begin{riskBox}{面试中的总原则}
当前项目是单用户自托管系统，但单用户不等于无认证，也不等于可以省略 owner scope。面试时应把“已实现”“已测试”“设计预留”和“尚未完成”分开说。尤其是生产密钥、HTTPS、真实外部模型和多主体部署，不能用本地单元测试替代现场证据。
\end{riskBox}

\subsection{可退化架构：从全有或全无变成健康、降级、恢复}

\begin{qaBox}{面试官可能问：为什么不让依赖失败就直接报错？}
因为研究系统的能力不是同质的。PostgreSQL、认证密钥和审计写入决定安全事实，失败时必须拒绝服务；RAG、重排序、浏览器和备用模型则是可选能力，失败时可以缩小研究范围并继续给出明确标注的结果。把所有依赖放进一个 try/except 会把安全依赖和可选能力混在一起，既可能扩大安全风险，也会把一次局部故障放大成整个应用不可用。
\end{qaBox}

\paragraph{设计。}
在 \term{app/security/capabilities.py} 中把外部能力抽象成 \term{CapabilityHealth} 状态机：\term{healthy} 表示最近探测成功，\term{degraded} 表示最近调用失败但应用有替代路径，\term{recovering} 表示冷却期后正在进行单飞探测。首次使用才初始化能力；失败记录 \term{reason\_code} 与时间；默认冷却 30 秒；并发请求共享一次恢复探测，避免所有请求同时下载模型或连接死服务。

降级结果不能只返回空列表，而应携带 \term{capability}、\term{state}、\term{reason\_code}、\term{retry\_after\_seconds} 和 \term{fallback\_used}。因此 DAG、SSE、审计和最终报告都能回答三个问题：哪个能力没用上，为什么没用上，什么时候可以再试。RAG 失败时跳过内部检索并标注“未使用内部知识库”；reranker 失败时保留向量/BM25 结果；浏览器失败时仍可保留搜索；报告流在首个正文输出前才允许切备用模型。

\paragraph{备选方案与取舍。}
方案 A 是任一依赖失败即终止。它的优点是语义简单、不会产生“部分结果”，缺点是可用性差、局部故障会放大，而且用户不知道已有证据是否仍然有价值。方案 B 是每个请求临时判断依赖是否可用。它比全局终止灵活，但并发下容易产生探测风暴，且不同节点可能对同一个能力做出不一致判断。方案 C 是集中式能力状态机加按能力定义的降级矩阵。项目选择方案 C，因为它同时解决状态一致性、节流恢复和用户可解释性；代价是增加状态维护、测试矩阵和错误码设计。

\paragraph{收益、代价与边界。}
收益是故障半径可控、恢复不需要重启、SSE 能继续表达进度，且安全关键依赖不会为了可用性而放行。代价是“部分完成”比成功/失败二元状态更难测试，也要求报告诚实标注证据缺口。已经考虑了冷却、单飞、可选能力回退和最终报告边界；尚未完成的是生产级指标告警、跨进程共享能力状态和真正外部模型的现场恢复演练。当前测试证明的是本进程状态机和降级语义，不是第三方服务的 SLO。

\subsection{安全边界：数据归属、令牌生命周期与默认拒绝}

\begin{qaBox}{面试官可能问：单用户系统为什么还要 owner scope？}
因为单用户是部署规模，不是数据访问模型。只要存在研究会话、文档、技能、配置、审批和审计，就需要明确它们属于哪个主体。严格的 owner 条件可以阻止历史数据被错误读取，也让未来增加第二主体时不必重写所有查询。
\end{qaBox}

\paragraph{设计。}
认证子系统使用 bcrypt 保存密码哈希，首次初始化通过 PostgreSQL advisory transaction lock 串行化，保留用户名拒绝被用于伪装系统主体。服务端 session 和 API token 只保存哈希值；明文 token 仅在创建响应中返回一次。每次校验同时检查撤销时间、过期时间、主体启用状态和 \term{session\_version}；密码轮换或统一撤销时递增版本并撤销旧 session/token，因此旧凭据不能继续访问。

所有资源查询使用主体条件，跨主体对象统一返回 404，避免泄露对象存在性。认证、审批和审计属于关键安全状态，不能因为 Redis 缓存故障而退化为放行。安全中间件统一设置精确 CORS、CSP、点击劫持防护、Referrer Policy、Content-Type 防嗅探和安全 Cookie；认证、上传和 token 路由还应使用独立限流桶。

\paragraph{备选方案与取舍。}
方案 A 只绑定 localhost，部署简单但无法安全远程访问。方案 B 完全信任反向代理，性能好但应用自身失去独立安全边界，代理配置错误会直接变成越权。方案 C 使用应用层密码、服务端 session 和可撤销 API token，并把 HTTPS 代理作为额外防线。项目采用方案 C，因为它可独立部署、能撤销泄露凭据，且不把安全责任转嫁给代理；代价是要维护认证表、锁、轮换和限流逻辑。

\paragraph{已经考虑与尚未考虑。}
已经考虑了 session fixation、密码轮换、删除/停用主体后的旧 token 失效、跨用户资源 404、Redis 故障不放行和 owner scope 回归。尚未在本轮引入多用户角色、组织/租户层级、WebAuthn、设备绑定和集中身份提供商；这些不能在面试中说成当前功能。

\subsection{配置与密钥：按主体隔离，并讲清威胁模型}

\begin{qaBox}{面试官可能问：应用层加密到底防什么？}
它主要防数据库文件、备份和日志泄露后直接暴露 LLM API Key。它不防已经获得进程权限的攻击者，因为运行中的进程必须能够解密并调用模型。把保护对象和不保护对象说清楚，比声称“密钥绝对安全”更可信。
\end{qaBox}

\paragraph{设计。}
LLM 配置属于受认证保护的管理员资源，数据库中的 API Key 使用部署主密钥加密，并以 \term{enc:} 前缀区分已迁移值。读取和写入均是幂等的，旧明文可以渐进迁移，不会因为重复启动而重复加密。配置 key 使用 \term{llm:<owner\_id>}，查询同时带 \term{owner\_id}，旧单管理员 \term{llm} 记录只在同一 owner 范围内兼容读取。生产模式拒绝缺失或示例主密钥、不安全 Cookie 和通配 CORS。

\paragraph{方案比较。}
明文入库最简单但数据库备份即等于泄露；只放环境变量能减少数据库暴露面，却失去运行时配置和管理员资源审计；应用层加密兼顾运行时配置、备份安全和渐进迁移，但引入主密钥轮换、启动失败和密钥丢失恢复问题。项目选择第三种，并明确“主密钥不与数据库备份放在一起”。

\paragraph{收益、代价与遗漏。}
收益是泄露数据库不直接泄露模型密钥，配置更新可审计且不必重训练。代价是主密钥轮换需要版本化密文和迁移工具，主密钥丢失会让密文不可恢复。已经考虑了前缀幂等迁移、数据库/备份威胁模型和生产启动检查；尚未完成 KMS/HSM 托管、自动轮换、密钥版本回滚和真实生产密钥演练。

\subsection{兼容层与 facade：重构不应一次性破坏调用面}

\term{app/agent\_tools} 的兼容层保留旧模块导入面，再逐步把解析、schema、执行和具体工具实现拆到更清晰的边界。它解决的是重构风险，不是业务功能本身：旧调用方继续拿到稳定 facade，新代码可以使用更细粒度的 adapter、policy 和执行器，迁移完成后再逐步收窄旧接口。

方案 A 是一次性重写所有调用方，结构最干净，但变更面大、回归定位困难；方案 B 是永远保留一层万能兼容代码，迁移成本低，却会形成双重语义和死代码；方案 C 是带弃用策略、契约测试和观测的渐进 facade。项目采用 C。收益是降低重构破坏性、支持分阶段发布和回滚；代价是短期内需要维护两套入口、增加文档和测试成本。已经考虑了旧导入兼容和 schema/执行边界；尚未完成的是统一弃用时间表、调用方迁移完成度指标和删除旧 facade 的最终版本。

\subsection{统一 LLM adapter：把 provider 差异隔离在边界}

\paragraph{设计。}
业务 Agent 不直接判断 provider，而通过 endpoint resolver 与 adapter 获取规范化的 base URL、认证头、模型名和客户端。OpenAI 兼容端点、Ollama 与 Anthropic 的差异集中在 \term{app/endpoint\_resolver.py} 和 \term{app/llm\_adapters.py}。HTTP 客户端共享连接池，并分别配置 connect/read/write/pool timeout；死端点进入冷却，有限重试只对可恢复的网络错误生效。

非流式调用可以在首个响应前切备用模型；流式调用把“是否已经向用户输出正文”作为协议状态。一旦已有正文，主模型失败只能发出可恢复的终止错误，不能无提示地切到备用模型，否则用户会看到语义断裂的半份报告。usage 缺失时使用保守估算，并在审计中标记 \term{estimated}。

\paragraph{为什么不是业务层 if/else。}
把 provider 判断散落在 Agent 中，短期代码少，长期会让 URL、认证头、超时、模型过滤和 fallback 规则彼此漂移。统一 adapter 的优点是替换模型和增加 provider 的影响面小、协议测试集中；缺点是 adapter 本身成为高复杂度边界，错误归类和原生 provider 特性可能被过度抽象。当前已覆盖 URL、认证头、模型筛选、usage 估算和首字节规则；尚未覆盖所有供应商的真实流式协议、工具调用差异和跨区域熔断指标。

\subsection{提示注入：把不可信数据放进数据流，而不是一句 system prompt}

网页、邮件、上传文档、RAG 片段、Skill 外部文本以及 MCP/tool 输出都必须携带来源、信任级别和内容边界。它们只能作为明确标注的用户上下文，不能拼接进 system role，也不能把其中的“请执行某命令”解释为控制指令。\term{prompt\_security} 负责数据封装，guardrail 负责研究政策，两者要分开建模：前者解决“这段文本能不能影响指令解释”，后者解决“这个问题是否允许联网、写入或输出”。

收益是攻击者无法仅靠网页文本升级为系统指令，且审计可以追踪污染来源；代价是 prompt 更长、模型可能降低可读性，且标记本身不是形式化证明。已经考虑了常见外部文本边界和 system role 隔离；尚未完成跨模型的注入基准、工具参数级 taint tracking、图像/多模态载荷和红队样本库。

\subsection{健康检查：liveness 与 readiness 必须分开}

\term{/health} 只回答进程是否活着，不能因为数据库暂时不可用就让容器被重启风暴拖垮。\term{/ready} 回答“能不能接流量”：数据库可连、schema 可用、认证主密钥有效、审计可写时才返回 200；RAG、Redis、浏览器和备用模型以不泄露敏感细节的信息状态返回。Docker healthcheck 应指向 \term{/ready}，因为编排系统需要依据接流量能力摘除实例。

方案 A 只有一个 health endpoint，简单但无法区分进程存活与业务可用；方案 B 只检查 HTTP server，误把“能响应”当成“能安全服务”；方案 C 分离 liveness/readiness 并分级关键依赖。项目采用 C，代价是部署和测试需要理解两种状态，且 readiness 不能泄露连接串、密钥或内部异常堆栈。当前代码已覆盖数据库、Redis 和认证状态的基本展示；真实容器冷启动和代理转发仍需现场验证。

\subsection{上传与文档生命周期：不是一个 save(file)}

上传链路应按顺序执行认证、owner 校验、文件名与路径规范化、扩展名/MIME/魔数交叉校验、大小/页数/解压限制、文本提取超时、哈希去重、主体配额和按主体限流。成功入库后，source 与 chunks 使用同一 owner；删除 source 时同步删除其 chunks 和向量记录；检索查询默认带 owner 条件。不要把用户原始路径写入宿主机，也不要用文件扩展名代替内容检测。

收益是降低路径逃逸、伪造 MIME、压缩炸弹、重复存储和跨主体检索风险；代价是解析库、缩略图、EXIF 处理和异步清理会增加复杂度，严格限制也可能拒绝合法的大文件。已经考虑了路径、MIME、哈希、owner、限流和同步删除；尚未完成的重点是对每一种文档格式做系统 fuzzing、病毒扫描、对象存储生命周期和可恢复的异步清理队列。

\subsection{备份导出/导入：merge-only 仍然要有事务边界}

导出只包含当前管理员拥有的认可资源，不导出密钥、session、token 和审计秘密。导入先做版本号与白名单结构校验，再按内容哈希去重、补齐 owner，并在同一个数据库连接的单一 transaction 中完成整批写入。任一 source 失败，整个批次回滚；成功响应同时返回 imported 与 skipped，保证幂等和冲突可见。

覆盖写入看起来简单，但会丢失当前数据、破坏 owner 关系，失败中途还很难恢复；逐条提交虽然易实现，却会留下半批导入。merge-only + 单事务的代价是大批次长事务占用连接、锁和 WAL，需要进一步做批量上限、分段导入协议或 staging table。当前已实现版本化、白名单、哈希 merge、owner 绑定和整批回滚测试；尚未完成跨版本迁移、加密备份容器、断点续传和大规模导入压测。

\subsection{本地持久化：JSON 也要遵守数据库式原子写}

仍需写本地文件的配置、Harness 状态和 session 快照，应统一使用同目录临时文件、完整写入、\term{flush}、\term{fsync} 和 \term{os.replace}。同目录保证替换发生在同一文件系统，\term{os.replace} 保证读者不会看到半截 JSON。写入工具应统一处理权限、编码、临时文件清理和 Windows 行为，禁止模块各自随手 \term{json.dump()}。

原子写不能解决磁盘损坏、错误数据逻辑或跨文件事务；它只保证单文件替换边界。因此 Harness 与数据库双写仍需分域 source of truth 和补偿恢复。收益是断电、进程终止和并发读时不容易得到损坏文件；代价是多一次磁盘写入、fsync 延迟和磁盘空间峰值。当前应把这套规则作为所有新增本地持久化的硬约束，而不能声称它等价于完整数据库事务。

\subsection{测试文化：每个真实事故都留下可回归的名字}

测试不应只验证 happy path，而应围绕曾经发生或可能发生的故障命名：初始化并发、错误密码、旧 token 失效、跨主体读取、短暂依赖故障、30 秒冷却、单飞恢复、首字节后禁止切换、恶意上传、伪造 MIME、导入回滚、提示注入和 readiness 拒流量。测试名本身就是维护文档，告诉未来维护者为什么一个看似多余的分支必须保留。

覆盖率只是信号，不是安全证明。90\% 以上可以降低未执行路径风险，但不能证明 threat model 完整；必须结合 property-based/fuzz 测试、迁移升级测试、真实协议测试和红队样本。当前自动化验证覆盖安全/adapter 94\%、全仓测试 172 项通过；尚未完成的是生产依赖现场、容器冷启动、多进程并发和长期故障注入。

\subsection{不在本轮范围的后台自动化与未来多主体}

定时任务、事件总线、邮件、CalDAV、跨设备同步和复杂角色体系不属于本轮生产重构边界。面试时可以说明其设计借鉴点：调度器要处理时区、短月钳制、幂等、单飞、串行执行和重启触发保留；事件必须按 owner 路由，不能一个用户的事件触发所有人的任务。但当前项目没有把这些能力伪装成已完成功能。保留 owner schema 和审计边界，是为未来扩展提供兼容接口，而不是提前承诺多用户产品。

\subsection{最终取舍总表}

\begin{center}
\small
\begin{tabularx}{\textwidth}{>{\bfseries}lXXX}
\toprule
主题 & 采用方案 & 主要收益 & 主要未覆盖项 \\
\midrule
依赖故障 & 能力状态机 + 分级降级 & 可用性与安全边界分离 & 跨进程状态、真实 SLO \\
访问控制 & 应用层认证 + owner scope & 可撤销、可独立部署 & 多角色、WebAuthn \\
密钥 & 应用层加密 + 渐进迁移 & 降低数据库泄露影响 & KMS/HSM、自动轮换 \\
LLM & 统一 adapter + 首字节协议 & provider 可替换、流式不串语义 & 全 provider 现场协议 \\
上传 & 生命周期校验 + 哈希 + owner & 减少文件攻击和越权 & fuzz、病毒扫描、对象存储 \\
备份 & 白名单 merge + 单批事务 & 幂等且不破坏现有数据 & 大规模断点续传 \\
健康 & liveness/readiness 分离 & 编排语义准确 & 生产代理现场演练 \\
测试 & 事故驱动回归 + 覆盖率 & 故障可复现、边界可审计 & 长期故障注入 \\
\bottomrule
\end{tabularx}
\end{center}

\begin{qaBox}{两分钟收口答案}
我没有把所有外部依赖都当成同一种错误，而是先按安全关键程度分级：数据库、认证密钥和审计写入失败时拒绝服务，RAG、浏览器、重排序和模型端点失败时通过能力状态机进入降级并持续探测恢复。所有资源查询坚持 owner scope，token 只存哈希且支持版本化撤销；API Key 用部署主密钥加密，但我会明确它防的是数据库/备份泄露，不防已攻陷进程。LLM provider 差异集中在 adapter，流式响应一旦输出正文就不允许静默切模型。上传、备份、提示注入和健康检查都有明确边界。最后我会把已测试能力、设计预留和尚未完成的生产现场演练分开说，而不是用覆盖率或文档替代真实验证。
\end{qaBox}


\section{附录：当前仓库现实、边界与演进路线}

\begin{summaryBox}{附录怎么用}
附录不是给面试官从头念的，而是给你自己在面试前校准边界。只要这里写的是“未完成”，面试时就不要包装成“已经具备”。只要这里写的是“设计已定”，面试时就要明确补一句“但代码尚未完整接入”。
\end{summaryBox}

\subsection{当前已核实的现实结论}
\begin{itemize}[leftmargin=2em]
  \item 聚合器覆盖失败语义的历史 bug 已经修复，不应继续当作现存问题对外表述。
  \item \term{tool\_call\_audit} 已落地，工具审计不再只依赖 session JSON。
  \item \term{session\_budget\_state} 与 \term{session\_runtime\_state} 已落地，预算和运行时状态已从收尾汇总前移到运行中持续更新。
  \item session 级 breaker 已有可用硬熔断版，但 token/cost 精度仍未完全做满。
  \item 动作级审批已经接入独立 reviewer 分支、审批表和审批 API。
  \item Harness supervisor 已接入最小可用实现，支持图外任务快照和批次边界恢复。
  \item MCP policy proxy 与 registry 已接入，但默认业务图还未大规模依赖 MCP 节点。
  \item skill 体系已经落地最小可用闭环：支持 Markdown + YAML frontmatter 解析、数据库持久化、CRUD、启动加载、热更新 registry、session 自动匹配与手动启用/禁用。
  \item skill 已不只是 prompt 附件，而是会真实影响 Planner、Analyst / Reflection / Report 提示，以及 session 级工具 allowlist。
\end{itemize}

\subsection{当前仍存在的关键边界}
\begin{riskBox}{需要诚实面对的点}
第一，审批虽然已经进入动作级执行流，但恢复粒度主要还是批次边界。第二，tool audit 和 budget state 已经 crash-safe 强化，但还不是全量法证平台。第三，Harness supervisor 已落地最小版本，但还没形成成熟的多 worker 调度和租约竞争体系。第四，MCP proxy 已存在，但默认 Planner 仍不以 MCP 为主要节点来源。第五，计费系统仍受上游 usage/billing 可观测性限制，不能吹成绝对精确。第六，skill 当前主要影响提示和工具边界，还没有演进到“自定义执行器 / 自定义 DAG 节点类型”这一层。
\end{riskBox}

\subsection{建议的三阶段演进}
\begin{enumerate}[leftmargin=2em]
  \item Phase 1：失败语义显式化、tool audit 持久化、会话级 budget breaker、session budget state。
  \item Phase 2：把最小可用 Harness 演进成多 worker supervisor，并把批次边界恢复做得更稳；同时把 skill 从“提示和工具约束层”继续推进到更精细的 session 策略解释和 API 治理。
  \item Phase 3：扩 Planner 的 MCP DAG 生成策略、细化 per-tool capability、接入外部审批系统信号，并评估 skill 是否需要升级成可挂载的自定义节点模板体系。
\end{enumerate}

\subsection{最终的面试总结句}
\begin{qaBox}{收口}
如果只看主链路，这个项目已经是一个能工作的研究型 Agent。真正让我认为它有工程价值的地方，在于它没有停留在“会规划、会搜索、会写报告”这一步，而是开始把失败语义、预算、审计、审批和外部工具边界都纳入系统设计。也正因为如此，它既能讲 Agent 能力，也能讲 Agent 治理。
\end{qaBox}

\end{document}
